\chapter{\'{E}quations canoniques}

La formulation des lois de la M\'{e}canique \`{a} l'aide de la fonction de Lagrange suppose, pour d\'{e}crire l'\'{e}tat m\'{e}canique d'un syst\`{e}me, les coordonn\'{e}es et les vitesses g\'{e}n\'{e}ralis\'{e}es de celui-ci. Sa description \`{a} l'aide des coordonn\'{e}es et des impulsions g\'{e}n\'{e}ralis\'{e}es pr\'{e}sente des avantages, notamment pour l'\'{e}tude de probl\`{e}mes g\'{e}n\'{e}raux de M\'{e}canique. Il reste donc \`{a} d\'{e}terminer les \'{e}quations du mouvement correspondant \`{a} cette m\'{e}thode.

En math\'{e}matique, la transformation de Legendre permet de passer d'un choix de variables ind\'{e}pendantes \`{a} un autre. Nous \'{e}tudierons cela ci-dessous.

La fonction de Lagrange est fonction des couples $(\{q_{i}\},\{\dot{q}_{i})$) donc :
\benn
	\mathrm{d}L = \sum_{i}\frac{\partial L}{\partial q_{i}}\mathrm{d}q_{i} + \sum_{i}\frac{\partial L}{\partial \dot{q}_{i}}\mathrm{d}\dot{q}_{i}
\eenn
Par application des formules (\ref{EQ:7_5}) et (\ref{EQ:7_7}) pour l'impulsion, nous pouvons \'{e}crire :
\be
	\mathrm{d}L = \sum_{i}\dot{p}_{i}\mathrm{d}q_{i} + \sum_{i}p_{i}\mathrm{d}\dot{q}_{i} \label{EQ:40_1}
\ee
or :
\benn
	\mathrm{d}\left(\sum_{i}p_{i}\dot{q}_{i}\right) = \sum_{i}\mathrm{d}p_{i}\dot{q}_{i} + \sum_{i}p_{i}\mathrm{d}\dot{q}_{i} \Leftrightarrow \sum_{i}p_{i}\mathrm{d}\dot{q}_{i} = \mathrm{d}\left(\sum_{i}p_{i}\dot{q}_{i}\right) - \sum_{i}\dot{q}_{i}\mathrm{d}p_{i}
\eenn
Ainsi, nous obtenons :
\bea
	\mathrm{d}L & = & \sum_{i}\dot{p}_{i}\mathrm{d}q_{i} + \mathrm{d}\left(\sum_{i}p_{i}\dot{q}_{i}\right) - \sum_{i}\dot{q}_{i}\mathrm{d}p_{i} \nonumber \\
	\Leftrightarrow \mathrm{d}\left(\sum_{i}\dot{p}_{i}\mathrm{d}q_{i} - L\right) & = & \sum_{i}\dot{q}_{i}\mathrm{d}p_{i} - \sum_{i}\dot{p}_{i}\mathrm{d}q_{i} \nonumber
\eea
Or $\sum_{i}\dot{p}_{i}\mathrm{d}q_{i} - L$ est la d\'{e}finition de l'\'{e}nergie totale $E$ du syst\`{e}me comme pr\'{e}sent\'{e} au paragraphe (\ref{PAR:6}). Son expression en fonction des coordonn\'{e}es et des impulsions est appel\'{e}e \emph{fonction de Hamilton} du syst\`{e}me et s'exprime ainsi :
\be
	H(q,p,t) = \sum_{i}p_{i}\dot{q}_{i} - L \label{EQ:40_2}
\ee
qui ob\'{e}it \`{a} :
\be
	\mathrm{d}H = -\sum_{i}\dot{p}_{i}\mathrm{d}q_{i} + \sum_{i}\dot{q}_{i}\mathrm{d}p_{i} \label{EQ:40_3}
\ee
pour laquelle les variables ind\'{e}pendantes sont les coordonn\'{e}es et les impulsions g\'{e}n\'{e}ralis\'{e}es. Par d\'{e}finition :
\benn
	\mathrm{d}H = \sum_{i}\frac{\partial H}{\partial q_{i}}\mathrm{d}q_{i} + \sum_{i}\frac{\partial H}{\partial p_{i}}\mathrm{d}p_{i}
\eenn
ce qui permet d'en d\'{e}duire :
\be
	\forall i\text{, }\dot{q}_{i} = \sum_{i}\frac{\partial H}{\partial p_{i}}\text{ et }\dot{p}_{i} = -\sum_{i}\frac{\partial H}{\partial q_{i}} \label{EQ:40_4}
\ee
Ces deux ensembles d'\'{e}quations sont les \'{e}quations du mouvement recherch\'{e}es et sont appel\'{e}es les \emph{\'{e}quations de Hamilton}. Elles constituent pour un syst\`{e}me \`{a} $s$ degr\'{e}s de libert\'{e}, $2s$ \'{e}quations diff\'{e}rentielles du premier ordre \`{a} $2s$ fonctions inconnues : $q(\mathrm{t})$ et $p(\mathrm{t})$. Elles remplacent les $s$ \'{e}quations du second ordre obtenues par la m\'{e}thode de Lagrange, voir les relations (\ref{EQ:2_6}). Les \'{e}quations de Hamilton sont dites \emph{canoniques} de par leur simplicit\'{e} et leur sym\'{e}trie.

La d\'{e}riv\'{e}e totale par rapport au temps de la fonction de Hamilton vaut :
\benn
	\frac{\mathrm{d}H}{\mathrm{dt}} = \frac{\partial H}{\partial\mathrm{t}} + \sum_{i}\frac{\partial H}{\partial q_{i}}\frac{\mathrm{d}q_{i}}{\mathrm{dt}} + \sum_{i}\frac{\partial H}{\partial p_{i}}\frac{\mathrm{d}p_{i}}{\mathrm{dt}} = \frac{\partial H}{\partial\mathrm{t}} + \sum_{i}\frac{\partial H}{\partial q_{i}}\dot{q}_{i} + \sum_{i}\frac{\partial H}{\partial p_{i}}\dot{p}_{i}
\eenn
Connaissant les \'{e}quations (\ref{EQ:40_4}), nous avons alors :
\be
	\frac{\mathrm{d}H}{\mathrm{dt}} = \frac{\partial H}{\partial\mathrm{t}} + \sum_{i}\frac{\partial H}{\partial q_{i}}\frac{\partial H}{\partial p_{i}} - \sum_{i}\frac{\partial H}{\partial p_{i}}\frac{\partial H}{\partial q_{i}} = \frac{\partial H}{\partial\mathrm{t}} \label{EQ:40_5}
\ee
Si la fonction de Hamilton ne d\'{e}pend pas explicitement du temps alors $\frac{\partial H}{\partial\mathrm{t}} = 0$ et nous retrouvons la loi de conservation de l'\'{e}nergie vue au paragraphe (\ref{PAR:6}).

En plus des couples $(q,\dot{q})$ pour la fonction de Lagrange et $(q,p)$ pour la fonction de Hamilton, ces deux fonctions peuvent \'{e}galement contenir des param\`{e}tres autres caract\'{e}risant le syst\`{e}me : param\`{e}tres intrins\`{e}ques ou champ ext\'{e}rieur\footnote{Le c{\oe}fficient de frottement ou l'amplitude de la force ext\'{e}rieure en sont deux exemples d\'{e}j\`{a} \'{e}voqu\'{e}s dans ce document.}. Soit $\lambda$ un de ces param\`{e}tres que nous consid\'{e}rons comme variable. La relation (\ref{EQ:40_1}) s'\'{e}crit alors :
\benn
	\mathrm{d}L = \sum_{i}\dot{p}_{i}\mathrm{d}q_{i} + \sum_{i}p_{i}\mathrm{d}\dot{q}_{i} + \frac{\partial L}{\partial\lambda}\mathrm{d}\lambda
\eenn
et (\ref{EQ:40_3}) :
\benn
	\mathrm{d}H = -\sum_{i}\dot{q}_{i}\mathrm{d}q_{i} + \sum_{i}p_{i}\mathrm{d}\dot{q}_{i} + \frac{\partial H}{\partial\lambda}\mathrm{d}\lambda
\eenn
Donc en reprenant la d\'{e}finition (\ref{EQ:40_2}) de $H$ :
\bea
	\mathrm{d}H & = & \sum_{i}p_{i}\mathrm{d}\dot{q}_{i} + \sum_{i}\dot{q}_{i}\mathrm{d}p_{i} - \mathrm{d}L \Leftrightarrow \left(\frac{\partial H}{\partial\lambda}\right)_{q,p}\mathrm{d}\lambda = \left(\frac{\partial L}{\partial\lambda}\right)_{q,\dot{q}}\mathrm{d}\lambda \nonumber \\
	\Leftrightarrow \left(\frac{\partial H}{\partial\lambda}\right)_{q,p} & = & \left(\frac{\partial L}{\partial\lambda}\right)_{q,\dot{q}} \label{EQ:40_6}
\eea
sachant que les quantit\'{e}s $\left(\frac{\partial H}{\partial\lambda}\right)_{q,p}$ est calcul\'{e}e pour $q$ et $p$ contants et $\left(\frac{\partial L}{\partial\lambda}\right)_{q,\dot{q}}$ pour $q$ et $\dot{q}$ constants.

De la m\^{e}me mani\`{e}re, si $L = L_{0} + L'$ est la fonction de Lagrange d'un syst\`{e}me telle que $L'$ est une quantit\'{e} petite compar\'{e}e \`{a} la fonction fondamentale $L_{0}$. En utilisant (\ref{EQ:40_2}) :
\benn
	H = \sum_{i}p_{i}\dot{q}_{i} - L = \sum_{i}p_{i}\dot{q}_{i} - L_{0} - L' = H_{0} - L'
\eenn
Nous pouvons donc \'{e}crire avec la m\^{e}me notation que :
\be
	\left(H'\right)_{q,p} = -\left(L'\right)_{q,\dot{q}} \label{EQ:40_7}
\ee
Lors du passage des relations (\ref{EQ:40_1}) \`{a} (\ref{EQ:40_3}), le temps n'a pas \'{e}t\'{e} pris en compte. Il joue le r\^{o}le d'un param\`{e}tre comme pr\'{e}sent\'{e} dans la relation (\ref{EQ:40_6}) et nous avons de fait :
\be
	\left(\frac{\partial H}{\partial\mathrm{t}}\right)_{q,p} = \left(\frac{\partial L}{\partial\mathrm{t}}\right)_{q,\dot{q}} \label{EQ:40_8}
\ee

\section{Fonction de Routh}\label{PAR:41}

Dans certains cas, pour le passage du couple de coordonn\'{e}es $(q,\dot{q})$ au couple $(q,p)$, il peut \^{e}tre utile de ne pas remplacer toutes les vitesses g\'{e}n\'{e}ralis\'{e}es par les impulsions mais seulement pour quelques unes.

Pour simplifier l'\'{e}criture et la d\'{e}monstration, nous supposons d'abord que nous n'avons que deux coordonn\'{e}es $q$ et $\xi$ et nous effectuons la transformation des variables $(q,\xi,\dot{q},\dot{\xi})$ vers les variables $(q,\xi,p,\dot{\xi})$ avec $p$ l'impulsion g\'{e}n\'{e}ralis\'{e}e correspondant \`{a} la variable $q$.

Le diff\'{e}rentiel de la fonction de Lagrange dans ce cas s'exprime ainsi :
\bea
	\mathrm{d} L & = & \frac{\partial L}{\partial q}\mathrm{d}q + \frac{\partial L}{\partial \dot{q}}\mathrm{d}\dot{q} + \frac{\partial L}{\partial \xi}\mathrm{d}\xi + \frac{\partial L}{\partial \dot{q}}\mathrm{d}\dot{\xi} \nonumber \\
	& = & \dot{p}\mathrm{d}q + p\mathrm{d}\dot{q} + \frac{\partial L}{\partial \xi}\mathrm{d}\xi + \frac{\partial L}{\partial \dot{q}}\mathrm{d}\dot{\xi} \nonumber
\eea
Et :
\benn
	\mathrm{d}(L - p\dot{q}) = \mathrm{d}L - \dot{q}\mathrm{d}p - p\mathrm{d}\dot{q} = \dot{p}\mathrm{d}q - \dot{q}\mathrm{d}p + \frac{\partial L}{\partial \xi}\mathrm{d}\xi + \frac{\partial L}{\partial \dot{q}}\mathrm{d}\dot{\xi}
\eenn
Nous introduisons d\'{e}sormais la \emph{fonction de Routh} :
\be
	R(q,p,\xi,\dot{\xi}) = p\dot{q} - L \label{EQ:41_1}
\ee
Par diff\'{e}renciation, nous \'{e}crivons :
\bea
	\mathrm{d}R & = & p\mathrm{d}\dot{q} + \dot{q}\mathrm{d}p - \dot{p}\mathrm{d}q - p\mathrm{d}\dot{q} - \frac{\partial L}{\partial \xi}\mathrm{d}\xi - \frac{\partial L}{\partial \dot{q}}\mathrm{d}\dot{\xi} \nonumber \\
	\Leftrightarrow \mathrm{d}R & = & -\dot{p}\mathrm{d}q + \dot{q}\mathrm{d}p - \frac{\partial L}{\partial \xi}\mathrm{d}\xi - \frac{\partial L}{\partial \dot{q}}\mathrm{d}\dot{\xi} \label{EQ:41_2}
\eea
et comme par d\'{e}finition de $R$, nous avons :
\benn
	\mathrm{d}R = \frac{\partial R}{\partial q}\mathrm{d}q + \frac{\partial L}{\partial \dot{q}}\mathrm{d}\dot{q} + \frac{\partial L}{\partial \xi}\mathrm{d}\xi + \frac{\partial L}{\partial \dot{q}}\mathrm{d}\dot{\xi}
\eenn
cela permet, par identification, de conclure \`{a} :
\bea
	\dot{q} = \frac{\partial R}{\partial p} & \text{ et } & \dot{p} = -\frac{\partial R}{\partial q} \label{EQ:41_3} \\
	\frac{\partial R}{\partial\xi} = -\frac{\partial L}{\partial\xi} & \text{ et } & \frac{\partial R}{\partial\dot{\xi}} = -\frac{\partial L}{\partial\dot{\xi}} \label{EQ:41_4}
\eea
Comme l'\'{e}quation de Lagrange pour la coordonn\'{e}e $\xi$ et d\'{e}finie par l'\'{e}quation (\ref{EQ:2_6}) est :
\be
	\frac{\mathrm{d}}{\mathrm{dt}}\left(\frac{\partial L}{\partial\dot{\xi}}\right) = \frac{\partial L}{\partial\xi} \Rightarrow \frac{\mathrm{d}}{\mathrm{dt}}\left(\frac{\partial R}{\partial\dot{\xi}}\right) = \frac{\partial R}{\partial\xi} \label{EQ:41_5}
\ee
Par analogie avec les \'{e}quations du mouvement exprim\'{e}es avec la fonction de Hamilton (\ref{EQ:40_4}), la fonction de Routh fait office de fonction de Hamilton pour les variables $q$ et $p$ et par analogie entre les \'{e}quations (\ref{EQ:2_6}) et (\ref{EQ:41_5}), la fonction de Routh fait office de fonction de Lagrange pour la coordonn\'{e}e $\xi$.

Par application de la d\'{e}finition de l'\'{e}nergie (\ref{EQ:6_1}), nous pouvons d\'{e}velopper :
\bea
	E & = & \dot{q}\frac{\partial L}{\partial\dot{q}} + \dot{\xi}\frac{\partial L}{\partial\dot{\xi}} - L = p\dot{q} + \dot{\xi}\frac{\partial L}{\partial\dot{\xi}} - L \nonumber \\
	\Leftrightarrow E & = & R + \dot{\xi}\frac{\partial L}{\partial\dot{\xi}} \Leftrightarrow E = R - \dot{\xi}\frac{\partial R}{\partial\dot{\xi}} \label{EQ:41_6}
\eea
Par lin\'{e}arit\'{e} des \'{e}quations pr\'{e}c\'{e}dentes, l'extension des formules obtenues pour le cas de plusieurs coordonn\'{e}es $q$ et $\xi$ est triviale.

La fonction de Routh est particuli\`{e}rement utile lorsque nous avons des coordonn\'{e}es cycliques, i.e. telles que si $\lambda$ est cyclique alors $\frac{\partial L}{\partial\lambda} = 0$ et la fonction de Lagrange $L$ ne d\'{e}pend donc pas explicitement de la variable $\lambda$.

Si les coordonn\'{e}es $q$ sont cycliques alors la fonction de Lagrange, de facto \'{e}galement la fonction de Routh, ne d\'{e}pendent que de $p$, $\xi$ et $\dot{\xi}$. L'\'{e}quation du mouvement pour la coordonn\'{e}e cyclique $q$ est :
\benn
	\frac{\mathrm{d}}{\mathrm{dt}}\left(\frac{\partial L}{\partial\dot{q}}\right) = \frac{\partial L}{\partial q} = 0 \Rightarrow \frac{\partial L}{\partial\dot{q}} = p = \text{cste}
\eenn
Ainsi, les impulsions associ\'{e}es aux coordonn\'{e}es cycliques sont constantes. Les \'{e}quations (\ref{EQ:41_5}) se r\'{e}duisent \`{a} des \'{e}quations contenant seulement les coordonn\'{e}es $\xi$. Si ces \'{e}quations sont r\'{e}solues et les fonctions $\xi(\mathrm{t})$ trouv\'{e}es alors la relation (\ref{EQ:41_3}) pour $\dot{q}$ permet de donner :
\benn
	\dot{q} = \frac{\partial R(p,\xi,\dot{\xi})}{\partial p}
\eenn
permettant ainsi de d\'{e}terminer $q(\mathrm{t})$ par int\'{e}gration directe.

\section{Crochets de Poisson}\label{PAR:42}

Nous consid\'{e}rons $f(q,p,\mathrm{t})$ une fonction des coordonn\'{e}es, des impulsions et du temps. Sa d\'{e}riv\'{e}e totale par rapport au temps s'\'{e}crit :
\bea
	\frac{\mathrm{d}f}{\mathrm{dt}} & = & \frac{\partial f}{\partial\mathrm{t}} + \sum_{k}\frac{\partial f}{\partial q_{k}}\dot{q}_{k} + \sum_{k}\frac{\partial f}{\partial p_{k}}\dot{p}_{k} \nonumber \\
	& = & \frac{\partial f}{\partial\mathrm{t}} + \sum_{k}\left(\frac{\partial f}{\partial q_{k}}\dot{q}_{k} + \frac{\partial f}{\partial p_{k}}\dot{p}_{k}\right) \nonumber
\eea
En utilisant les \'{e}quations de Hamilton en (\ref{EQ:40_4}), on a :
\benn
	\frac{\mathrm{d}f}{\mathrm{dt}} = \frac{\partial f}{\partial\mathrm{t}} + \sum_{k}\left(\frac{\partial H}{\partial p_{k}}\frac{\partial f}{\partial q_{k}} - \frac{\partial H}{\partial q_{k}}\frac{\partial f}{\partial p_{k}}\right)
\eenn
que nous pouvons \'{e}crire ainsi :
\be
	\frac{\mathrm{d}f}{\mathrm{dt}} = \frac{\partial f}{\partial\mathrm{t}} + \{H,j\} \label{EQ:42_1}
\ee
en introduisant les \emph{crochets de Poisson} pour $H$ et $f$ tels que :
\be
	\{H,f\} = \sum_{k}\left(\frac{\partial H}{\partial p_{k}}\frac{\partial f}{\partial q_{k}} - \frac{\partial H}{\partial q_{k}}\frac{\partial f}{\partial p_{k}}\right) \label{EQ:42_2}
\ee
De telles fonctions des variables dynamiques qui restent constantes dans le mouvement constituent des intégrales premi\`{e}res, voir le paragraphe (\ref{PAR:6}). Dans le cas (\ref{EQ:42_1}), $f$ est int\'{e}grale premi\`{e}re du mouvement si et seulement si :
\be
	\frac{\mathrm{d}f}{\mathrm{dt}} = 0 \Leftrightarrow \frac{\partial f}{\partial\mathrm{t}} + \{H,f\} = 0 \label{EQ:42_3}
\ee
ou encore si $f$ ne d\'{e}pend pas explicitement du temps :
\be
	\{H,f\} = 0 \label{EQ:42_4}
\ee
De mani\`{e}re triviale, les crochets de Poisson de deux grandeurs quelconques vallent :
\be
	\{f,g\} = \sum_{k}\left(\frac{\partial f}{\partial p_{k}}\frac{\partial g}{\partial q_{k}} - \frac{\partial f}{\partial q_{k}}\frac{\partial g}{\partial p_{k}}\right) \label{EQ:42_5}
\ee
D\'{e}terminons d\'{e}sormais quelques propri\'{e}t\'{e}s :
\benn
        \{f,g\} = \sum_{k}\left(\frac{\partial f}{\partial p_{k}}\frac{\partial g}{\partial q_{k}} - \frac{\partial f}{\partial q_{k}}\frac{\partial g}{\partial p_{k}}\right) = - \sum_{k}\left(\frac{\partial f}{\partial q_{k}}\frac{\partial g}{\partial p_{k}} - \frac{\partial f}{\partial p_{k}}\frac{\partial g}{\partial q_{k}}\right)
\eenn
soit :
\be
	\{f,g\} = -\{g,f\} \label{EQ:42_6}
\ee
Si $c$ est une constante alors :
\be
	\forall k\text{, } \frac{\partial c}{\partial q_{k}} = \frac{\partial c}{\partial p_{k}} = 0 \Rightarrow \{f,c\} = 0 \label{EQ:42_7}
\ee
Par lin\'{e}arit\'{e} de la d\'{e}rivation :
\be
	\{f_{1} + f_{2},g\} = \{f_{1},g\} + \{f_{2},g\} \label{EQ:42_8}
\ee
De plus :
\bea
	\{f_{1}f_{2},g\} & = & \sum_{k}\left(\frac{\partial f_{1}f_{2}}{\partial p_{k}}\frac{\partial g}{\partial q_{k}} - \frac{\partial f_{1}f_{2}}{\partial q_{k}}\frac{\partial g}{\partial p_{k}}\right) \nonumber \\
	& = & \sum_{k}\left(f_{2}\frac{\partial f_{1}}{\partial p_{k}}\frac{\partial g}{\partial q_{k}} + f_{1}\frac{\partial f_{2}}{\partial p_{k}}\frac{\partial g}{\partial q_{k}} - f_{2}\frac{\partial f_{1}}{\partial q_{k}}\frac{\partial g}{\partial p_{k}} - f_{1}\frac{\partial f_{2}}{\partial q_{k}}\frac{\partial g}{\partial p_{k}}\right) \nonumber \\
	& = & f_{1}\sum_{k}\left(\frac{\partial f_{2}}{\partial p_{k}}\frac{\partial g}{\partial q_{k}} - \frac{\partial f_{2}}{\partial q_{k}}\frac{\partial g}{\partial p_{k}}\right) + f_{2}\sum_{k}\left(\frac{\partial f_{1}}{\partial p_{k}}\frac{\partial g}{\partial q_{k}} - \frac{\partial f_{1}}{\partial q_{k}}\frac{\partial g}{\partial p_{k}}\right) \nonumber \\
	\Leftrightarrow \{f_{1}f_{2},g\} & = & f_{1}\{f_{2},g\} + f_{2}\{f_{1},g\} \label{EQ:42_9}
\eea
De plus, en appliquant la d\'{e}riv\'{e}e partielle par rapport au temps :
\bea
	\frac{\partial\{f,g\}}{\partial\mathrm{t}} & = & \sum_{k}\left(\frac{\partial}{\partial\mathrm{t}}\left(\frac{\partial f}{\partial p_{k}}\right)\frac{\partial g}{\partial q_{k}} + \frac{\partial f}{\partial p_{k}}\frac{\partial}{\partial\mathrm{t}}\left(\frac{\partial g}{\partial q_{k}}\right) - \frac{\partial}{\partial\mathrm{t}}\left(\frac{\partial f}{\partial q_{k}}\right)\frac{\partial g}{\partial p_{k}} - \frac{\partial f}{\partial q_{k}}\frac{\partial}{\partial\mathrm{t}}\left(\frac{\partial g}{\partial p_{k}}\right)\right) \nonumber \\
	& = & \sum_{k}\left(\frac{\partial}{\partial p_{k}}\left(\frac{\partial f}{\partial\mathrm{t}}\right)\frac{\partial g}{\partial q_{k}} - \frac{\partial}{\partial q_{k}}\left(\frac{\partial f}{\partial\mathrm{t}}\right)\frac{\partial g}{\partial p_{k}} + \frac{\partial f}{\partial p_{k}}\frac{\partial}{\partial q_{k}}\left(\frac{\partial g}{\partial\mathrm{t}}\right) - \frac{\partial f}{\partial q_{k}}\frac{\partial}{\partial p_{k}}\left(\frac{\partial g}{\partial\mathrm{t}}\right)\right) \nonumber \\
	\Leftrightarrow \frac{\partial\{f,g\}}{\partial\mathrm{t}} & = & \left\{\frac{\partial f}{\partial\mathrm{t}},g\right\} + \left\{f,\frac{\partial g}{\partial\mathrm{t}}\right\} \label{EQ:42_10}
\eea
Prenons d\'{e}sormais quelques cas particuliers en remarquant que $\forall i\text{, }\frac{\partial q_{i}}{\partial p_{i}} = 0$ et $\forall i\ne k\text{, }\frac{\partial q_{i}}{\partial q_{k}} = 0$ :
\begin{itemize}
	\item si $g = q_{k}$ alors :
\be
	\{f,q_{k}\} = \frac{\partial f}{\partial p_{k}}\frac{\partial q_{k}}{\partial q_{k}} - 0 \Leftrightarrow \{f,q_{k}\} = \frac{\partial f}{\partial p_{k}} \label{EQ:42_11}
\ee
	\item si $g = p_{k}$ alors :
\be
	\{f,p_{k}\} = 0 - \frac{\partial f}{\partial q_{k}}\frac{\partial p_{k}}{\partial p_{k}} \Leftrightarrow \{f,p_{k}\} = -\frac{\partial f}{\partial q_{k}} \label{EQ:42_12}
\ee
	\item par extension :
\be
	\begin{cases}
		\forall i,k\text{, }\{q_{i},q_{k}\} = \frac{\partial q_{i}}{\partial p_{k}} = 0 \nonumber \\
		\forall i,k\text{, }\{p_{i},p_{k}\} = -\frac{\partial p_{i}}{\partial p_{k}} = 0 \nonumber \\
		\forall i,k\text{, }\{p_{i},q_{k}\} = \frac{\partial p_{i}}{\partial p_{k}} = \delta_{ik} \label{EQ:42:13}
	\end{cases}
\ee
	\item de plus, dans le cas o\`{u} $f$ est la fonction de Hamilton d'un syst\`{e}me alors :
\be
	\begin{cases}
		\forall k\text{, }\{H,q_{k}\} = \frac{\partial H}{\partial p_{k}} = \dot{q}_{k} \nonumber \\
		\forall k\text{, }\{H,p_{k}\} = -\frac{\partial H}{\partial q_{k}} = \dot{p}_{k} \nonumber
	\end{cases}
\ee
\end{itemize}
Nous allons maintenant \'{e}tudier la relation entre les crochets de Poisson form\'{e}s par trois fonctions, \`{a} savoir : $\{f,\{g,h\}\} + \{g,\{h,f\}\} + \{h,\{g,f\}\}$ dans le cas g\'{e}n\'{e}ral d'un syst\`{e}me compos\'{e} de $n$ particules.
Tout d'abord, calculons cette expression dans le cas particulier d'une unique particule. Le premier terme se d\'{e}veloppe ainsi :
\bea
	\{f,\{g,h\}\}_{n=0} & = & \frac{\partial f}{\partial p_{0}}\frac{\partial}{\partial q_{0}}\left(\frac{\partial g}{\partial p_{0}}\frac{\partial h}{\partial q_{0}} - \frac{\partial g}{\partial q_{0}}\frac{\partial h}{\partial p_{0}}\right) - \frac{\partial f}{\partial q_{0}}\frac{\partial}{\partial p_{0}}\left(\frac{\partial g}{\partial p_{0}}\frac{\partial h}{\partial q_{0}} - \frac{\partial g}{\partial q_{0}}\frac{\partial h}{\partial p_{0}}\right) \nonumber \\
	& = & \frac{\partial f}{\partial p_{0}}\left(\frac{\partial^{2}g}{\partial q_{0}\partial p_{0}}\frac{\partial h}{\partial q_{0}} + \frac{\partial g}{\partial p_{0}}\frac{\partial^{2}h}{\partial q_{0}^{2}} - \frac{\partial^{2}g}{\partial q_{0}^{2}}\frac{\partial h}{\partial p_{0}} - \frac{\partial g}{\partial q_{0}}\frac{\partial^{2}h}{\partial q_{0}p_{0}}\right) \nonumber \\
	& & -\frac{\partial f}{\partial q_{0}}\left(\frac{\partial^{2}g}{\partial p_{0}^{2}}\frac{\partial h}{\partial q_{0}} + \frac{\partial g}{\partial p_{0}}\frac{\partial^{2}h}{\partial q_{0}p_{0}} - \frac{\partial^{2}g}{\partial q_{0}p_{0}}\frac{\partial h}{\partial p_{0}} - \frac{\partial g}{\partial q_{0}}\frac{\partial^{2}h}{\partial p_{0}^{2}}\right) \nonumber
\eea
Le deuxi\`{e}me terme se d\'{e}veloppe ainsi :
\bea
	\{g,\{h,f\}\}_{n=0} & = & \frac{\partial g}{\partial p_{0}}\frac{\partial}{\partial q_{0}}\left(\frac{\partial h}{\partial p_{0}}\frac{\partial f}{\partial q_{0}} - \frac{\partial h}{\partial q_{0}}\frac{\partial f}{\partial p_{0}}\right) - \frac{\partial g}{\partial q_{0}}\frac{\partial}{\partial p_{0}}\left(\frac{\partial h}{\partial p_{0}}\frac{\partial f}{\partial q_{0}} - \frac{\partial h}{\partial q_{0}}\frac{\partial f}{\partial p_{0}}\right) \nonumber \\
        & = & \frac{\partial g}{\partial p_{0}}\left(\frac{\partial^{2}h}{\partial q_{0}\partial p_{0}}\frac{\partial f}{\partial q_{0}} + \frac{\partial h}{\partial p_{0}}\frac{\partial^{2}f}{\partial q_{0}^{2}} - \frac{\partial^{2}h}{\partial q_{0}^{2}}\frac{\partial f}{\partial p_{0}} - \frac{\partial h}{\partial q_{0}}\frac{\partial^{2}f}{\partial q_{0}p_{0}}\right) \nonumber \\
        & & -\frac{\partial g}{\partial q_{0}}\left(\frac{\partial^{2}h}{\partial p_{0}^{2}}\frac{\partial f}{\partial q_{0}} + \frac{\partial h}{\partial p_{0}}\frac{\partial^{2}f}{\partial q_{0}p_{0}} - \frac{\partial^{2}h}{\partial q_{0}p_{0}}\frac{\partial f}{\partial p_{0}} - \frac{\partial h}{\partial q_{0}}\frac{\partial^{2}f}{\partial p_{0}^{2}}\right) \nonumber
\eea
Et enfin le troisi\`{e}me crochet :
\bea
        \{h,\{f,g\}\}_{n=0} & = & \frac{\partial h}{\partial p_{0}}\frac{\partial}{\partial q_{0}}\left(\frac{\partial f}{\partial p_{0}}\frac{\partial g}{\partial q_{0}} - \frac{\partial f}{\partial q_{0}}\frac{\partial g}{\partial p_{0}}\right) - \frac{\partial h}{\partial q_{0}}\frac{\partial}{\partial p_{0}}\left(\frac{\partial f}{\partial p_{0}}\frac{\partial g}{\partial q_{0}} - \frac{\partial f}{\partial q_{0}}\frac{\partial g}{\partial p_{0}}\right) \nonumber \\
        & = & \frac{\partial h}{\partial p_{0}}\left(\frac{\partial^{2}f}{\partial q_{0}\partial p_{0}}\frac{\partial g}{\partial q_{0}} + \frac{\partial f}{\partial p_{0}}\frac{\partial^{2}g}{\partial q_{0}^{2}} - \frac{\partial^{2}f}{\partial q_{0}^{2}}\frac{\partial g}{\partial p_{0}} - \frac{\partial f}{\partial q_{0}}\frac{\partial^{2}g}{\partial q_{0}p_{0}}\right) \nonumber \\
        & & -\frac{\partial h}{\partial q_{0}}\left(\frac{\partial^{2}f}{\partial p_{0}^{2}}\frac{\partial g}{\partial q_{0}} + \frac{\partial f}{\partial p_{0}}\frac{\partial^{2}g}{\partial q_{0}p_{0}} - \frac{\partial^{2}f}{\partial q_{0}p_{0}}\frac{\partial g}{\partial p_{0}} - \frac{\partial f}{\partial q_{0}}\frac{\partial^{2}g}{\partial p_{0}^{2}}\right) \nonumber
\eea
et nous remarquons que pour la somme des trois \'{e}l\'{e}ments, tous les termes s'annulent deux \`{a} deux et :
\benn
	\{f,\{g,h\}\}_{n=0} = 0
\eenn
Cette identit\'{e} est valable pour n'importe quel syst\`{e}me n'ayant qu'une unique particule.
En proc\'{e}dant de la m\^{e}me mani\`{e}re, le calcul de la diff\'{e}rence $\{f,\{g,h\}\}_{n+1} - \{f,\{g,h\}\}_{n}$ donne la valeur du crochet de Poisson pour la particule identifi\'{e}e $n+1$ qui est nulle, comme d\'{e}montr\'{e} pr\'{e}c\'{e}demment. Nous pouvons donc en conclure que :
\be
	\{f,\{g,h\}\} + \{g,\{h,f\}\} + \{h,\{g,f\}\} = 0 \label{EQ:42_14}
\ee
qui est \emph{l'identit\'{e} de Jacobi}.

Une propri\'{e}t\'{e} importante des crochets de Poisson consiste en ce que, si $f$ et $g$ sont deux int\'{e}grales premi\`{e}res du mouvement alors leurs crochets de Poisson l'est aussi. C'est le \emph{th\'{e}or\`{e}me de Poisson}. Pour le d\'{e}motrer, nous commen\c{c}ons par appliquer l'identit\'{e} de Jacobi (\ref{EQ:42_14}) en identifiant $h$ \`{a} la fonction de Hamilton $H$, soit :
\benn
	\{f,\{g,H\}\} + \{g,\{H,f\}\} + \{H,\{g,f\}\} = 0
\eenn
La d\'{e}finition (\ref{EQ:42_1}) donne :
\benn
	\frac{\mathrm{d}}{\mathrm{dt}}\left(\{f,g\}\right) = \frac{\partial}{\partial\mathrm{t}}\left(\{f,g\}\right) + \{H,\{f,g\}\} = \left\{\frac{\partial f}{\partial\mathrm{t}},g\right\} + \left\{f,\frac{\partial g}{\partial\mathrm{t}}\right\} + \{H,\{f,g\}\}
\eenn
par application de la formule (\ref{EQ:42_10}). L'identit\'{e} (\ref{EQ:42_14}) permet de d\'{e}velopper la relation pr\'{e}c\'{e}dente telle que :
\bea
	\frac{\mathrm{d}}{\mathrm{dt}}\left(\{f,g\}\right) & = & \left\{\frac{\partial f}{\partial\mathrm{t}},g\right\} + \left\{f,\frac{\partial g}{\partial\mathrm{t}}\right\} - \{f,\{g,H\}\} - \{g,\{H,f\}\} \nonumber \\
	& = & \left\{\frac{\partial f}{\partial\mathrm{t}},g\right\} + \left\{f,\frac{\partial g}{\partial\mathrm{t}}\right\} + \{\{H,g\},f\} + \{\{H,f\},g\} \nonumber \\
	& = & \left\{\frac{\partial f}{\partial\mathrm{t}} + \{H,f\},g\right\} + \left\{f,\frac{\partial g}{\partial\mathrm{t}} + \{H,g\}\right\} \nonumber \\
	\Leftrightarrow \frac{\mathrm{d}}{\mathrm{dt}}\left(\{f,g\}\right) & = & \left\{\frac{\mathrm{d}f}{\mathrm{dt}},g\right\} + \left\{f,\frac{\mathrm{d}g}{\mathrm{dt}}\right\} \label{EQ:42_16}
\eea
Or $f$ et $g$ sont des int\'{e}grales premi\`{e}res, i.e. $\frac{\mathrm{d}f}{\mathrm{dt}} = \frac{\mathrm{d}g}{\mathrm{dt}} = 0$ donc :
\benn
	\frac{\mathrm{d}}{\mathrm{dt}}\left(\{f,g\}\right) = 0
\eenn
$\{f,g\}$ est donc bien une int\'{e}grale premi\`{e}re si $f$ et $g$ le sont aussi. Ce th\'{e}or\`{e}me s'applique \'{e}galement si $f$ et $g$ ne d\'{e}pendent pas explicitement du temps.

Enfin et parce que dans un syst\`{e}me fini compos\'{e} de $s$ degr\'{e}s de libert\'{e}, le nombre maximal d'int\'{e}grales premi\`{e}res est de $2s - 1$, voir le paragraphe (\ref{PAR:6}), alors le th\'{e}or\`{e}me de Poisson ne permet donc pas toujours d'obtenir de nouvelles int\'{e}grales premi\`{e}res.

\section{L'action en fonction des coordonn\'{e}es}\label{PAR:43}

L'action d\'{e}finie dans la relation (\ref{EQ:2_1}) est l'int\'{e}grale de la fonction de Lagrange du syst\`{e}me m\'{e}canique :
\be
	S = \int_{t_{1}}^{t_{2}}L\mathrm{dt} \label{EQ:43_1}
\ee
prise le long de la trajectoire entre deux positions donn\'{e}es $q^{(1)}$ et $q^{(2)}$ occup\'{e}es par le syst\`{e}me aux instants $\mathrm{t}_{1}$ et $\mathrm{t}_{2}$. La comparaison des valeurs de l'int\'{e}grale pour des trajectoires voisines ayant les m\^{e}mes limites $q(\mathrm{t}_{1})$ et $q(\mathrm{t}_{2})$ fait varier l'action. Une seule de ces trajectoires correspond au mouvement r\'{e}el, c'est celle pour laquelle l'action $S$ est minimum.

\subsection{Cas o\`{u} seule la coordonn\'{e}e finale est variable}

Nous allons d\'{e}sormais d\'{e}terminer l'action $S$ comme une grandeur caract\'{e}risant le mouvement le long de trajectoires r\'{e}elles. Nous comparerons alors les valeurs de $S$ pour des trajectoires ayant une origine commune $q_{1} = q(\mathrm{t}_{1})$ mais passant \`{a} l'instant $\mathrm{t}_{2}$ par des positions diff\'{e}rentes. En d'autres termes, nous consid\'{e}rerons l'int\'{e}grale d'action pour des trajectoires r\'{e}elles comme fonction des coordonn\'{e}es \`{a} la limite sup\'{e}rieure d'int\'{e}gration.

La variation de l'action lorsque nous passons d'une trajectoire \`{a} une trajectoire voisine est donn\'{e}e pour un seul degr\'{e} de libert\'{e} par la relation (\ref{EQ:2_5}) :
\benn
	\delta S = \left[\frac{\partial L}{\partial\dot{q}}\delta q\right]_{\mathrm{t}_{1}}^{\mathrm{t}_{2}} - \int_{\mathrm{t}_{1}}^{\mathrm{t}_{2}}{\left(\frac{\partial L}{\partial q} - \frac{\mathrm{d}}{\mathrm{dt}}\left(\frac{\partial L}{\partial\dot{q}}\right)\right)\delta q\mathrm{dt}}
\eenn
Puisque les trajectoires r\'{e}elles satisfont aux \'{e}quations de Lagrange alors l'int\'{e}grale de la relation pr\'{e}c\'{e}dente s'annule. Posons d\'{e}sormais $\delta q(\mathrm{t}_{1}) = 0$ et $\delta q(\mathrm{t}_{2}) = \delta q$ et puisque l'\'{e}quation d\'{e}finissant l'impulsion donne $p = \frac{\partial L}{\partial\dot{q}}$, nous avons $\delta S = p\delta q$ qui s'\'{e}crit pour un nombre quelconque de degr\'{e}s de libert\'{e} :
\be
	\delta S = \sum_{i}p_{i}\delta q_{i} \label{EQ:43_2}
\ee
et par voie de cons\'{e}quence :
\be
	\forall i\text{, }\frac{\partial S}{\partial q_{i}} = p_{i} \label{EQ:43_3}
\ee

\subsection{Cas o\`{u} la dur\'{e}e du mouvement est variable}

De fa\c{c}on analogue, nous pouvons envisager l'action comme fonction explicite du temps. Dans ce cas, nous consid\'{e}rons les trajectoires partant \`{a} un instant donn\'{e} $\mathrm{t}_{1}$ d'une position $q^{(1)}$ pour aboutir \`{a} une position donn\'{e}e $q^{(2)}$ \`{a} des instants diff\'{e}rents $\mathrm{t}_{2} = \mathrm{t}$. La quantit\'{e} $\frac{\partial S}{\partial\mathrm{t}}$ peut \^{e}tre d\'{e}termin\'{e}e au moyen d'une variation appropri\'{e}e de l'int\'{e}grale de la relation (\ref{EQ:2_5}).

Par d\'{e}finition de l'action, voir la formule (\ref{EQ:2_1}), sa d\'{e}riv\'{e}e totale par rapport au temps le long des trajectoires vaut :
\be
	\frac{\mathrm{d}S}{\mathrm{dt}} = L \label{EQ:43_4}
\ee
Or :
\benn
	\frac{\mathrm{d}S}{\mathrm{dt}} = \frac{\partial S}{\partial\mathrm{t}} + \sum_{i}\frac{\partial S}{\partial q_{i}}\frac{\mathrm{d}q_{i}}{\mathrm{dt}}
\eenn
puisque, ici, $S$ n'est fonction que des coordonn\'{e}es. En appliquant la relation (\ref{EQ:43_3}), nous \'{e}crivons :
\bea
	\frac{\mathrm{d}S}{\mathrm{dt}} & = & \frac{\partial S}{\partial\mathrm{t}} + \sum_{i}p_{i}\dot{q}_{i} \Leftrightarrow \frac{\partial S}{\partial\mathrm{t}} = L - \sum_{i}p_{i}\dot{q}_{i} \nonumber \\
	\Leftrightarrow \frac{\partial S}{\partial\mathrm{t}} & = & -H \label{EQ:43_5}
\eea
qui est en accord avec la d\'{e}finition de la fonction de Hamilton $H$ de la formule (\ref{EQ:40_2}).

En prenant l'\'{e}galit\'{e} (\ref{EQ:43_3}), nous obtenons finalement :
\bea
	\frac{\mathrm{d}S}{\mathrm{dt}} & = & -H + \sum_{i}p_{i}\dot{q}_{i} = -H + \sum_{i}p_{i}\frac{\mathrm{d}q_{i}}{\mathrm{dt}} \nonumber \\
	\Leftrightarrow \mathrm{d}S & = & \sum_{i}p_{i}\mathrm{d}q_{i} - H\mathrm{dt} \label{EQ:43_6}
\eea
qui donne la diff\'{e}rentielle totale de l'action comme fonction des coordonn\'{e}es et du temps \`{a} la limite sup\'{e}rieure de l'int\'{e}grale (\ref{EQ:43_1}).

\subsection{Cas o\`{u} les coordonn\'{e}es initiale et finale et le temps sont variables}

Nous nous mettons dans le cadre o\`{u} les coordonn\'{e}es et le temps varient au d\'{e}but et \`{a} la fin du mouvement. La variation de l'action est donn\'{e}e par la diff\'{e}rence aux deux extremit\'{e}s de la relation (\ref{EQ:43_6}), soit :
\be
	\mathrm{d}S = \sum_{i}p_{i}^{(2)}\mathrm{d}q_{i}^{(2)} - H^{(2)}\mathrm{dt}^{(2)} - \sum_{i}p_{i}^{(1)}\mathrm{d}q_{i}^{(1)} + H^{(1)}\mathrm{dt}^{(1)} \label{EQ:43_7}
\ee
Quelle que soit l'action ext\'{e}rieure qui pourrait s'exercer sur le syst\`{e}me pendant son mouvement, l'\'{e}tat final du syst\`{e}me n'est pas une fonction arbitraire de son \'{e}tat initial. Seules sont possibles les trajectoires pour lesquelles la relation (\ref{EQ:43_7}) est une diff\'{e}rentielle totale exacte.

De plus, le principe de moindre action, ind\'{e}pendamment de la fonction de Lagrange, impose \`{a} l'ensemble des trajectoires possibles des limites bien d\'{e}termin\'{e}es.

\subsection{D\'{e}duction des \'{e}quations de Hamilton}

L'int\'{e}grale (\ref{EQ:43_6}) permet d'\'{e}crire :
\be
	S = \int{\left(\sum_{i}p_{i}\mathrm{d}q_{i} - H\mathrm{dt}\right)} \label{EQ:43_8}
\ee
qui suppose que les coordonn\'{e}es et les impulsions sont des grandeurs ind\'{e}pendantes. Nous simplifions par la suite avec une coordonn\'{e}e et une impulsion :
\benn
	S = \int{\left(p\mathrm{d}q - H\mathrm{dt}\right)}
\eenn
et sa variation devient :
\benn
	\delta S = \int{\left(\delta p\mathrm{d}q + p\mathrm{d}(\delta q) - \delta H\mathrm{dt}\right)}
\eenn
car $\mathrm{d}(\delta\mathrm{t}) = 0$. Or :
\benn
	\delta H = \frac{\partial H}{\partial q}\delta q + \frac{\partial H}{\partial p}\delta p
\eenn
qui permet d'\'{e}crire la variation de l'action telle que :
\bea
	\delta S & = & \int{\left(\delta p\mathrm{d}q + p\delta(\mathrm{d}q) - \frac{\partial H}{\partial q}\delta q\mathrm{dt} - \frac{\partial H}{\partial p}\delta p\mathrm{dt} \right)} \nonumber \\
	& = & \int{\left(\mathrm{d}q - \frac{\partial H}{\partial p}\mathrm{dt}\right)\delta p} + \int p\mathrm{d}(\delta q) - \int{\frac{\partial H}{\partial q}\delta q\mathrm{dt}} \nonumber
\eea
Une int\'{e}gration par partie de la quantit\'{e} $\int p\mathrm{d}(\delta q)$ donne :
\benn
	\int p\mathrm{d}(\delta q) = \left[p\delta q\right] - \int\delta q\mathrm{d}p
\eenn
que nous pouvons injecter dans la variation de l'action et obtenir :
\benn
	\delta S = \int{\delta p\left(\mathrm{d}q - \frac{\partial H}{\partial p}\mathrm{dt}\right)} + \left[p\delta q\right] - \int{\delta q\left(\mathrm{d}p - \frac{\partial H}{\partial q}\mathrm{dt}\right)}
\eenn
Nous posons $\left[p\delta q\right] = 0$ par choix des conditions aux limites. Et par application du principe de moindre action (\ref{EQ:2_4}), \`{a} savoir $\delta S = 0$, et puisque les quantit\'{e} $\delta q$ et $\delta p$ sont ind\'{e}pendantes et arbitraires, alors il devient n\'{e}cessaire de poser :
\bea
	\mathrm{d}q - \frac{\partial H}{\partial p}\mathrm{dt} = 0 & \text{ et } & \mathrm{d}p - \frac{\partial H}{\partial q}\mathrm{dt} = 0 \nonumber \\
	\frac{\partial H}{\partial p} = \dot{q} & \text{ et } & \frac{\partial H}{\partial q} = -\dot{p} \nonumber
\eea
Nous retrouvons donc bien les \'{e}quations de Hamilton d\'{e}finies par les relations (\ref{EQ:40_4}).

\section{Principe de Maupertuis}\label{PAR:44}

Le mouvement d'un syst\`{e}me m\'{e}canique est enti\`{e}rement d\'{e}termin\'{e} par le principe de moindre action. En r\'{e}solvant les \'{e}quations du mouvement qui en d\'{e}coulent, nous d\'{e}terminons la forme de la trajectoire ainsi que la relation entre la position et le temps \'{e}coul\'{e}.

Nous nous concentrons maintenant uniquement sur la question plus restreinte de la d\'{e}termination de la forme de la trajectoire. Dans notre cas pr\'{e}cis, les fonctions de Lagrange et de Hamilton ne d\'{e}pendent pas explicitement du temps et par cons\'{e}quent, l'\'{e}nergie du syst\`{e}me se conserve : $H(p,q) = E$.

Le principe de moindre action stipule que pour des valeurs initiale et finale des coordonn\'{e}es et du temps donn\'{e}es, la variation de l'action est nulle. Mais si nous autorisons la variation de l'instant final $\mathrm{t}$, les coordonn\'{e}es initiale et finale restant toujours fixes, alors la relation (\ref{EQ:43_5}) s'applique :
\be
	\delta S = -H\delta\mathrm{t} \label{EQ:44_1}
\ee
Dans le cadre de la conservation de l'\'{e}nergie, la relation pr\'{e}c\'{e}dente (\ref{EQ:44_1}) devient :
\be
	\delta S + E\delta\mathrm{t} = 0 \label{EQ:44_2}
\ee
et la relation (\ref{EQ:43_8}) devient : $S = \int(\sum_{i}p_{i}\mathrm{d}q_{i} - H\mathrm{dt}) = \int\sum_{i}p_{i}\mathrm{d}q_{i} - E\int\mathrm{dt}$ ou encore :
\be
	S = \int\sum_{i}p_{i}\mathrm{d}q_{i} - E(\mathrm{t} - \mathrm{t}_{0}) \label{EQ:44_3}
\ee
\`{A} partir de maintenant, nous d\'{e}finissons l'\emph{action r\'{e}duite} telle que :
\be
	S_{0} = \int\sum_{i}p_{i}\mathrm{d}q_{i} \label{EQ:44_4}
\ee
et la formule (\ref{EQ:44_3}) se r\'{e}\'{e}crit :
\benn
	\delta S = \delta S_{0} - E\delta\mathrm{t}
\eenn
Or, en accord avec la relation (\ref{EQ:44_2}), nous pouvons en d\'{e}duire :
\be
	\delta S_{0} = 0 \label{EQ:44_5}
\ee
Ainsi l'action r\'{e}duite a un minimum sur l'ensemble des trajectoires satisfaisant \`{a} la loi de conservation de l'\'{e}nergie et passant par le point final \`{a} un instant arbitraire.

Pour utiliser le principe variationnel exprim\'{e} en (\ref{EQ:44_5}), il nous faut d'abord exprimer les impulsions en fonction des coordonn\'{e}es $q$ et de leurs diff\'{e}rentielles $\mathrm{d}q$. Ce principe variationnel est appel\'{e} \emph{principe de Maupertuis}.

Reprenons la formule de la fonction de Lagrange (\ref{EQ:5_5}) :
\benn
	L = \frac{1}{2}\sum_{i,j}a_{ij}(q)\dot{q}_{i}\dot{q}_{j} - U(q)
\eenn
avec les impulsions :
\bea
	p_{i} & = & \frac{\partial L}{\partial\dot{q}_{i}} = \frac{1}{2}\sum_{i,j}a_{ij}\frac{\partial\dot{q}_{i}\dot{q}_{j}}{\partial\dot{q}_{i}} = \frac{1}{2}\sum_{i=j}a_{ij}\frac{\partial\dot{q}_{i}^{2}}{\partial\dot{q}_{i}} \nonumber \\
	\Leftrightarrow p_{i} & = & \sum_{i,j}a_{ij}\dot{q}_{j} \nonumber
\eea
et l'\'{e}nergie constante :
\benn
	E = \frac{1}{2}\sum_{i,j}a_{ij}(q)\dot{q}_{i}\dot{q}_{j} + U(q) = \frac{1}{2}\sum_{i,j}a_{ij}(q)\frac{\mathrm{d}q_{i}}{\mathrm{dt}}\frac{\mathrm{d}q_{j}}{\mathrm{dt}} + U(q)
\eenn
qui permet d\'{e}crire la diff\'{e}rientelle en temps comme :
\be
	\mathrm{dt} = \sqrt{\frac{\sum_{i,j}a_{ij}\mathrm{d}q_{i}\mathrm{d}q_{j}}{2(E - U)}} \label{EQ:44_8}
\ee
Nous pouvons alors d\'{e}velopper la diff\'{e}rentielle de l'action r\'{e}duite :
\bea
	\mathrm{d}S_{0} & = & \sum_{i}p_{i}\mathrm{d}q_{i} = \sum_{i,j}a_{ij}\dot{q}_{j}\mathrm{d}q_{i} = \sum_{i,j}a_{ij}\frac{\mathrm{d}q_{j}}{\mathrm{dt}}\mathrm{d}q_{i} \nonumber \\
	& = & \sqrt{2(E - U)}\frac{\sum_{i,j}a_{ij}\mathrm{d}q_{j}\mathrm{d}q_{i}}{\sqrt{\sum_{i,j}a_{ij}\mathrm{d}q_{j}\mathrm{d}q_{i}}} = \sqrt{2(E - U)\sum_{i,j}a_{ij}\mathrm{d}q_{j}\mathrm{d}q_{i}} \nonumber
\eea
Soit, finalement :
\be
	S_{0} = \int{\sqrt{2(E - U)\sum_{i,j}a_{ij}\mathrm{d}q_{j}\mathrm{d}q_{i}}} \label{EQ:44_9}
\ee
Prenons l'exemple d'un point mat\'{e}riel de masse $m$. Son \'{e}nergie cin\'{e}tique est $T = \frac{m}{2}\left(\frac{\mathrm{d}l}{\mathrm{dt}}\right)^{2}$ avec $\mathrm{d}l$ un \'{e}l\'{e}ment de longueur le long de la trajectoire. Le principe variationnel donne :
\benn
	\delta S_{0} = \delta\left(\int{\sqrt{2(E - U)m(\mathrm{d}l)^{2}}}\right) = \delta\left(\int{\sqrt{2m(E - U)}\mathrm{d}l}\right)
\eenn
soit, par application de la formule (\ref{EQ:44_8}) :
\be
	\delta\left(\int{\sqrt{2m(E - U)}\mathrm{d}l}\right) = 0 \label{EQ:44_10}
\ee
o\`{u} l'int\'{e}grale est calcul\'{e}e entre deux points donn\'{e}es de l'espace.

Pour le mouvement libre d'une particule, soit tel que $U = 0$, l'\'{e}quation (\ref{EQ:44_10}) devient :
\benn
	\delta\left(\int{\sqrt{2mE}\mathrm{d}l}\right) = 0 \Leftrightarrow \delta\left(\int\mathrm{d}l\right) = 0
\eenn
car $E$ et $m$ sont des constantes par rapport au temps. Ainsi la particule suit le chemin le plus court, \`{a} savoir une droite.

Nous reprenons maintenant l'expression (\ref{EQ:44_3}) et nous op\'{e}rons une variation par rapport \`{a} l'\'{e}nergie $E$ :
\benn
	S = S_{0} - E(\mathrm{t} - \mathrm{t}_{0}) \Rightarrow \delta S = \frac{\partial S_{0}}{\partial E}\delta E - \delta E(\mathrm{t} - \mathrm{t}_{0}) - E\delta\mathrm{t}
\eenn
Or la relation (\ref{EQ:44_2}) donne $\delta S = -E\delta\mathrm{t}$, donc nous obtenons :
\be
	\frac{\partial S_{0}}{\partial E} = \mathrm{t} - \mathrm{t}_{0} \label{EQ:44_11}
\ee
En prenant l'action r\'{e}duite sous sa forme (\ref{EQ:44_9}), nous \'{e}crivons :
\bea
	\frac{\partial S_{0}}{\partial E} & = & \frac{\partial}{\partial E}\left(\int{\sqrt{2(E - U)\sum_{i,j}a_{ij}\mathrm{d}q_{j}\mathrm{d}q_{i}}}\right) = \int{\frac{\partial}{\partial E}\left(\sqrt{2(E - U)\sum_{i,j}a_{ij}\mathrm{d}q_{j}\mathrm{d}q_{i}}\right)} \nonumber \\
	& = & \int{\sqrt{\sum_{i,j}a_{ij}\mathrm{d}q_{j}\mathrm{d}q_{i}}\frac{\partial}{\partial E}\left(\sqrt{2(E - U)}\right)} = \int{\frac{\sum_{i,j}a_{ij}\mathrm{d}q_{j}\mathrm{d}q_{i}}{\sqrt{2(E - U)}}} \nonumber
\eea
Ainsi la relation (\ref{EQ:44_11}) devient dans ce cadre :
\be
	\int{\frac{\sum_{i,j}a_{ij}\mathrm{d}q_{j}\mathrm{d}q_{i}}{\sqrt{2(E - U)}}} = \mathrm{t} - \mathrm{t}_{0} \label{EQ:44_12}
\ee
qui n'est rien d'autre que l'int\'{e}grale de la relation (\ref{EQ:44_8}). Avec l'\'{e}quation de la trajectoire, l'int\'{e}grale d\'{e}termine enti\`{e}rement le mouvement.

\section{Transformations canoniques}\label{PAR:45}

\subsection{D\'{e}finitions}

Le choix des coordonn\'{e}es g\'{e}n\'{e}ralis\'{e}es $q$ n'est limit\'{e} par aucune condition et nous pouvons prendre $s$ grandeurs quelconques d\'{e}finissant de fa\c{c}on univoque la position du syst\`{e}me \'{e}tudi\'{e} dans l'espace. Les \'{e}quations de Lagrange (\ref{EQ:2_6}) ne d\'{e}pendent pas de ce choix et nous pouvons en d\'{e}duire qu'elles sont invariantes par rapport \`{a} la transformation qui fait passer des coordonn\'{e}es $q_{1}$, $q_{2}$, $\ldots$, $q_{s}$ \`{a} d'autres grandeurs ind\'{e}pendantes $Q_{1}$, $Q_{2}$, $\ldots$, $Q_{s}$. Les nouvelles coordonn\'{e}es $Q_{i}$ sont des fonctions de $q_{i}$. Admettons que cette relation contienne \'{e}galement le temps, i.e. :
\be
	\forall i\text{, }Q_{i} = Q_{i}(q,\mathrm{t}) \label{EQ:45_1}
\ee
qui sont $s$ \emph{transformations ponctuelles}.

De la m\^{e}me mani\`{e}re que pour les \'{e}quations de Lagrange, la transformation (\ref{EQ:45_1}) laisse invariantes les \'{e}quations de Hamilton (\ref{EQ:40_4}).

Toutefois, dans la m\'{e}thode de Hamilton, les impulsions $p$ jouent le r\^{o}le de variables ind\'{e}pendantes au m\^{e}me titre que les coordonn\'{e}es $q$. Ainsi la notion de transformation peut \^{e}tre \'{e}largie pour englober la transformation des $2s$ variables ind\'{e}pendantes $p$ et $q$ en de nouvelles variables $P$ et $Q$ suivant les formules :
\be
	\forall i\text{, }Q_{i} = Q_{i}(p,q,\mathrm{t})\text{ et }P_{i} = P_{i}(p,q,\mathrm{t}) \label{EQ:45_2}
\ee
D\'{e}gageons les conditions auxquelles doit ob\'{e}ir une transformation pour que les \'{e}quations du mouvement avec les nouvelles variables $P$ et $Q$ aient la forme connue :
\be
	\dot{Q}_{i} = \frac{\partial H'}{\partial P_{i}}\text{ et }\dot{P}_{i} = -\frac{\partial H'}{\partial Q_{i}} \label{EQ:45_3}
\ee
avec $H'(P,Q)$ une nouvelle fonction de Hamilton. De telles transformations sont dites \emph{transformations canoniques}. Elles ne changent que la forme des \'{e}quations de Hamilton.

\subsection{Formules des transformations canoniques}\label{PAR:45_1}

\`{A} partir de la relation de l'action en fonction de la fonction de Hamilton (\ref{EQ:43_8}), nous avons retrouv\'{e} les \'{e}quations de Hamilton en utilisant le principe de moindre action tel que :
\be
	\delta\left(\int{\left(\sum_{i}p_{i}\mathrm{d}q_{i} - H\mathrm{dt}\right)}\right) = 0 \label{EQ:45_4}
\ee
Pour que $P$ et $Q$ satisfassent aussi aux \'{e}quations de Hamilton, elles doivent v\'{e}rifier le principe de moindre action :
\be
	\delta\left(\int{\left(\sum_{i}P_{i}\mathrm{d}Q_{i} - H'\mathrm{dt}\right)}\right) = 0 \label{EQ:45_5}
\ee
Les deux principes (\ref{EQ:45_4}) et (\ref{EQ:45_5}) ne sont \'{e}quivalents qu'\`{a} la condition que les expressions int\'{e}gr\'{e}es diff\`{e}rent par la diff\'{e}rentielle totale d'une fonction arbitraire $F$ des coordonn\'{e}es, des impulsions et du temps. La diff\'{e}rence des deux int\'{e}grales sera alors une constante dont la variation sera nulle, voir la relation (\ref{EQ:2_8}). Par cons\'{e}quent, nous pouvons en conclure :
\benn
	\sum_{i}p_{i}\mathrm{d}q_{i} - H\mathrm{dt} = \sum_{i}P_{i}\mathrm{d}Q_{i} - H'\mathrm{dt} + \mathrm{d}F
\eenn
Toute transformation canonique est caract\'{e}ris\'{e}e par sa fonction $F$ appel\'{e}e \emph{fonction g\'{e}n\'{e}ratrice de la transformation}. La relation devient alors :
\be
	\mathrm{d}F = \sum_{i}p_{i}\mathrm{d}q_{i} - \sum_{i}P_{i}\mathrm{d}Q_{i} + (H' - H)\mathrm{dt} \label{EQ:45_6}
\ee
Or par d\'{e}finition de la d\'{e}riv\'{e}e totale :
\benn
	\mathrm{d}F = \frac{\partial F}{\partial q_{i}}\mathrm{d}q_{i} + \frac{\partial F}{\partial p_{i}}\mathrm{d}p_{i} + \frac{\partial F}{\partial Q_{i}}\mathrm{d}Q_{i} + \frac{\partial F}{\partial P_{i}}\mathrm{d}P_{i} + \frac{\partial F}{\partial\mathrm{t}}\mathrm{dt}
\eenn
cela nous permet donc d'en d\'{e}duire :
\be
	p_{i} = \frac{\partial F}{\partial q_{i}}\text{, }P_{i} = \frac{\partial F}{\partial Q_{i}}\text{ et }H' = H + \frac{\partial F}{\partial\mathrm{t}} \label{EQ:45_7}
\ee
Nous voyons alors que la fonction g\'{e}n\'{e}ratrice est donn\'{e}e comme fonction des anciennes et des nouvelles coordonn\'{e}es et du temps : $F = F(q,Q,\mathrm{t})$.

Nous \'{e}tablissons maintenant la fonction g\'{e}n\'{e}ratrice en fonction des anciennes coordonn\'{e}es $q$ et des nouvelles impulsions $P$. La relation (\ref{EQ:45_6}) peut \'{e}galement s'\'{e}crire :
\bea
	\mathrm{d}F & = & \sum_{i}p_{i}\mathrm{d}q_{i} - \left(\mathrm{d}(\sum_{i}P_{i}Q_{i}) - \sum_{i}Q_{i}\mathrm{d}P_{i}\right) + (H' - H)\mathrm{dt} \nonumber \\
	\Leftrightarrow \mathrm{d}(F + \sum_{i}P_{i}Q_{i}) & = & \sum_{i}p_{i}\mathrm{d}q_{i} + \sum_{i}Q_{i}\mathrm{d}P_{i} + (H' - H)\mathrm{dt} \nonumber
\eea
La nouvelle fonction g\'{e}n\'{e}ratrice est $\Phi = F + \sum_{i}P_{i}Q_{i} = \Phi(q,P,\mathrm{t})$ avec :
\be
	p_{i} = \frac{\partial\Phi}{\partial q_{i}}\text{, }Q_{i} = \frac{\partial\Phi}{\partial P_{i}}\text{ et }H' = H + \frac{\partial \Phi}{\partial\mathrm{t}} \label{EQ:45_8}
\ee
Pour chaque transformation canonique, la diff\'{e}rence $H' - H$ est toujours \'{e}gale \`{a} la d\'{e}riv\'{e}e partielle par rapport au temps de la fonction g\'{e}n\'{e}ratrice. Dans le cas particulier o\`{u} cette derni\`{e}re ne d\'{e}pend pas du temps alors $H' = H$ et la nouvelle fonction de Hamilton s'obtient en portant dans $H$ les quantit\'{e}s $p$ et $q$ exprim\'{e}es en fonction des nouvelles variables $P$ et $Q$.

\subsection{Grandeurs canoniques conjugu\'{e}es}\label{PAR:45_2}

Avec la notion de transformation canonique dans la m\'{e}thode de Hamilton, les notions de coordonn\'{e}es et d'impulsions g\'{e}n\'{e}ralis\'{e}es perdent de leur sens physique initial. En effet, les transformations (\ref{EQ:45_2}) relient chaque grandeur $P$ et $Q$ aussi bien aux coordonn\'{e}es $q$ qu'aux impulsions $p$ et de fait, les variables ne peuvent plus \^{e}tre consid\'{e}r\'{e}es comme des coordonn\'{e}es purement spatiales. La diff\'{e}rence entre les deux groupes de variables est alors une question de nomenclature.

Elle peut se manifester de mani\`{e}re concr\`{e}te en prenant par exemple la transformation $Q_{i} = p_{i}$ et $P_{i} = -q_{i}$ pour laquelle la d\'{e}finition de la fonction g\'{e}n\'{e}ratrice donne :
\benn
	\mathrm{d}F = \sum_{i}Q_{i}\mathrm{d}q_{i} - \sum_{i}(-q_{i})\mathrm{d}Q_{i} + (H' - H)\mathrm{dt} = \mathrm{d}(\sum_{i}q_{i}\mathrm{d}Q_{i}) + (H' - H)\mathrm{dt}
\eenn
Cela ne change pas la forme canonique des \'{e}quations mais interchange les noms des coordonn\'{e}es et des variables. Compte tenu de la convention de terminologie, les variables $p$ et $q$ sont appel\'{e}es dans la m\'{e}thode de Hamilton des \emph{grandeurs canoniquement conjugu\'{e}es}.

La condition pour que des variables soient canoniquement conjugu\'{e}es peut \^{e}tre exprim\'{e}e avec des crochets de Poisson, voir le paragraphe (\ref{PAR:42}). Pour le d\'{e}montrer, nous commen\c{c}ons par le th\'{e}or\`{e}me d'invariance des crochets de Poisson aux transformations canoniques.

Soit $\{f,g\}_{p,q}$ le crochet de Poisson des grandeurs $f$ et $g$ dans lequel les d\'{e}rivations s'effectuent par rapport \`{a} $p$ et $q$, et $\{f,g\}_{P,Q}$ le crochet de Poisson des grandeurs $f$ et $g$ dans lequel les d\'{e}rivations s'effectuent par rapport \`{a} $P$ et $Q$.

Les transformations canoniques (\ref{EQ:45_7}) et (\ref{EQ:45_8}) montrent que le temps joue le r\^{o}le d'un param\`{e}tre et peut donc nous permettre de n'\'{e}tudier que le cas o\`{u} les grandeurs $f$ et $g$ ne d\'{e}pendent pas explicitement du temps. Nous assimilons d\'{e}sormais $g$ comme fonction de Hamilton d'un syst\`{e}me fictif. Alors :
\bea
	\{f,g\}_{p,q} & = & \sum_{i}\left(\frac{\partial f}{\partial p_{i}}\frac{\partial g}{\partial q_{i}} - \frac{\partial f}{\partial q_{i}}\frac{\partial g}{\partial p_{i}}\right) = \sum_{i}\left(\frac{\partial f}{\partial p_{i}}(-\dot{p}_{i}) - \frac{\partial f}{\partial q_{i}}\dot{q}_{i}\right) = -\sum_{i}\left(\frac{\partial f}{\partial p_{i}}(-\dot{p}_{i}) + \frac{\partial f}{\partial q_{i}}\dot{q}_{i}\right) \nonumber \\
	& = & -\frac{\mathrm{d}f}{\mathrm{dt}} \nonumber
\eea
en appliquant la relation (\ref{EQ:42_1}). Par d\'{e}finition, une d\'{e}rivation totale ne d\'{e}pend pas du choix des variables donc la quantit\'{e} $\frac{\mathrm{d}f}{\mathrm{dt}}$ est aussi \'{e}gale \`{a} $\{f,g\}_{P,Q}$. Cela nous permet donc d'en conclure que :
\be
	\{f,g\}_{p,q} = \{f,g\}_{P,Q} \label{EQ;45_9}
\ee
Par application directe, nous pouvons en conclure que :
\be
	\begin{cases}
		\{Q_{i},Q_{j}\}_{p,q} = \{Q_{i},Q_{j}\}_{P,Q} = \sum_{i}\left(\frac{\partial Q_{i}}{\partial P_{i}}\frac{\partial Q_{j}}{\partial Q_{i}} - \frac{\partial Q_{i}}{\partial Q_{i}}\frac{\partial Q_{j}}{\partial P_{i}}\right) = 0 \nonumber \\
		\{P_{i},P_{j}\}_{p,q} = \{P_{i},P_{j}\}_{P,Q} = \sum_{i}\left(\frac{\partial P_{i}}{\partial P_{i}}\frac{\partial P_{j}}{\partial Q_{i}} - \frac{\partial P_{i}}{\partial Q_{i}}\frac{\partial P_{j}}{\partial P_{i}}\right) = 0 \nonumber \\
		\{P_{i},Q_{j}\}_{p,q} = \{P_{i},Q_{j}\}_{P,Q} = \sum_{i}\left(\frac{\partial P_{i}}{\partial P_{i}}\frac{\partial Q_{j}}{\partial Q_{i}} - \frac{\partial P_{i}}{\partial Q_{i}}\frac{\partial Q_{j}}{\partial P_{i}}\right) = \delta_{ij} \label{EQ:45_10}
	\end{cases}
\ee
Ces trois relations sont les conditions auxquelles doivent satisfaire les nouvelles variables pour que la transformation $(p,q)\rightarrow(P,Q)$ soit canonique.

\subsection{Cas particulier du mouvement comme transformation canonique}\label{PAR:45_3}

Il est \`{a} noter que la variation des grandeurs $p$ et $q$ lors du mouvement peut \^{e}tre consid\'{e}r\'{e}e comme une transformation canonique.

Soit $p_{\mathrm{t}}$ et $q_{\mathrm{t}}$ les valeurs des variables canoniques \`{a} l'instant $\mathrm{t}$ et $p_{\mathrm{t+\tau}}$ et $q_{\mathrm{t+\tau}}$ leurs valeurs \`{a} l'instant $\mathrm{t} + \tau$. Ces derni\`{e}res sont fonction des premi\`{e}res et de la valeur $\tau$ de l'intervalle de temps en tant que param\`{e}tre :
\benn
	q_{\mathrm{t+\tau}} = q(q_{\mathrm{t}},p_{\mathrm{t}},\tau)\text{ et }p_{\mathrm{t+\tau}} = p(q_{\mathrm{t}},p_{\mathrm{t}},\tau)
\eenn
Si nous consid\'{e}rons ces formules comme une transformation faisant passer des variables $q_{\mathrm{t}},p_{\mathrm{t}}$ aux variables $q_{\mathrm{t+\tau}},p_{\mathrm{t+\tau}}$ alors cette transformation sera canonique. L'expression (\ref{EQ:43_7}) donne dans notre cas :
\bea
	\mathrm{S} & = & \sum_{i}p_{\mathrm{t}+\tau}\mathrm{d}q_{\mathrm{t}+\tau} - H'\tau - \sum_{i}p_{\mathrm{t}}\mathrm{d}q_{\mathrm{t}} - H\tau \nonumber \\
	& = & \sum_{i}p_{\mathrm{t}+\tau}\mathrm{d}q_{\mathrm{t}+\tau} - \sum_{i}p_{\mathrm{t}}\mathrm{d}q_{\mathrm{t}} - (H' - H)\tau \nonumber
\eea
relation identifiable \`{a} la d\'{e}finition de la fonction g\'{e}n\'{e}ratrice (\ref{EQ:45_6}) telle que :
\benn
	\mathrm{d}(-S) = \sum_{i}p_{\mathrm{t}}\mathrm{d}q_{\mathrm{t}} - \sum_{i}p_{\mathrm{t}+\tau}\mathrm{d}q_{\mathrm{t}+\tau} + (H' - H)\tau
\eenn
Alors $-S$ est la fonction g\'{e}n\'{e}ratrice de cette transformation canonique au cours du mouvement.

\section{Th\'{e}or\`{e}me de Liouville}\label{PAR:46}

Pour l'interpr\'{e}tation g\'{e}om\'{e}trique des ph\'{e}nom\`{e}nes m\'{e}caniques, nous utilisons la notion d'\emph{espace des phases}. C'est un espace \`{a} $2s$ dimensions sur les axes desquel sont port\'{e}es les $s$ coordonn\'{e}es g\'{e}n\'{e}ralis\'{e}es et les $s$ impulsions du syst\`{e}me. Chaque point de cet espace correspond \`{a} un \'{e}tat d\'{e}termin\'{e} du syst\`{e}me. Lorsque le syst\`{e}me se meut, le point de l'espace des phases qui le repr\'{e}sente d\'{e}crit une ligne appel\'{e}e trajectoire dans l'espace des phases.

Le produit des diff\'{e}rentielles $\mathrm{d}\Gamma = \mathrm{d}q_{1}\ldots\mathrm{d}q_{s}\mathrm{d}p_{1}\ldots\mathrm{d}p_{s}$ repr\'{e}sente un \'{e}l\'{e}ment de volume de l'espace des phases. Nous consid\'{e}rons $\int{\mathrm{d}\Gamma}$ pris sur un domaine quelconque de cet espace et repr\'{e}sentant le volume de ce domaine.

Nous allons d\'{e}montrer que cette grandeur poss\`{e}de la propri\'{e}t\'{e} d'\^{e}tre invariante par rapport aux transformations canoniques. Ainsi si nous transformons canoniquement $p,q$ en $P,Q$, nous aurons :
\be
	\int\cdots\int\mathrm{d}q_{1}\ldots\mathrm{d}q_{s}\mathrm{d}p_{1}\ldots\mathrm{d}p_{s} = \int\cdots\int\mathrm{d}Q_{1}\ldots\mathrm{d}Q_{s}\mathrm{d}P_{1}\ldots\mathrm{d}P_{s} \label{EQ:46_1}
\ee
La transformation des variables dans une int\'{e}grale multiple s'effectue en utilisant la valeur absolue du d\'{e}terminant jacobien\footnote{Voir \emph{https://fr.wikipedia.org/wiki/Matrice\textunderscore jacobienne} pour quelques g\'{e}n\'{e}ralit\'{e}s} telle que :
\benn
	\int\cdots\int\mathrm{d}Q_{1}\ldots\mathrm{d}Q_{s}\mathrm{d}P_{1}\ldots\mathrm{d}P_{s} = \int\cdots\int\lvert J\rvert\mathrm{d}q_{1}\ldots\mathrm{d}q_{s}\mathrm{d}p_{1}\ldots\mathrm{d}p_{s}
\eenn
avec $J$ le d\'{e}terminant de la matrice jacobienne d\'{e}finie ainsi :
\benn
	\{J_{ij}\} = 
	\begin{cases}
		\frac{\partial Q_{i}}{\partial q_{j}}\text{ pour }i\le s\text{ et }j\le s \\
		\frac{\partial P_{i}}{\partial q_{j}}\text{ pour }i > s\text{ et }j\le s \\
		\frac{\partial Q_{i}}{\partial p_{j}}\text{ pour }i\le s\text{ et }j > s \\
		\frac{\partial P_{i}}{\partial p_{j}}\text{ pour }i > s\text{ et }j > s
	\end{cases}
\eenn
qui est ordinairement not\'{e}e :
\benn
	\{J\} = \frac{\partial(Q_{1},\ldots,Q_{s},P_{1},\ldots,P_{s})}{\partial(q_{1},\ldots,q_{s},p_{1},\ldots,p_{s})}
\eenn
Nous devons alors d\'{e}montrer que dans le cadre d'une transformation canonique quelconque :
\be
	\lvert J\rvert = 1 \label{EQ:46_3}
\ee

En utilisant certaines r\`{e}gles des jacobiens, nous pouvons d\'{e}velopper tel que :
\bea
	\{J\} & = & \frac{\partial(Q_{1},\ldots,Q_{s},P_{1},\ldots,P_{s})}{\partial(q_{1},\ldots,q_{s},P_{1},\ldots,P_{s})} \bigg/ \frac{\partial(Q_{1},\ldots,Q_{s},P_{1},\ldots,P_{s})}{\partial(q_{1},\ldots,q_{s},P_{1},\ldots,P_{s})} \nonumber \\
	& = & \left\{\frac{(Q_{1},\ldots,Q_{s})}{q_{1},\ldots,q_{s}}\right\}_{P=cste} \bigg/ \left\{\frac{p_{1},\ldots,p_{s}}{P_{1},\ldots,P_{s}}\right\}_{q=cste} \label{EQ:46_4}
\eea
En \'{e}crivant les termes de la matrice jacobienne pour une fonction g\'{e}n\'{e}ratrice $\Phi$ et en s'appuyant sur les \'{e}galit\'{e}s (\ref{EQ:45_8}), nous pouvons \'{e}crire :
\benn
        \{J_{ij}\} =
        \begin{cases}
		\frac{\partial Q_{i}}{\partial q_{j}} = \frac{\partial^{2}\Phi}{\partial q_{j}\partial P_{i}}\text{ pour }i\le s\text{ et }j\le s \\
                \frac{\partial P_{i}}{\partial q_{j}} = \frac{\partial^{2}\Phi}{\partial q_{j}\partial q_{i}}\text{ pour }i > s\text{ et }j\le s \\
                \frac{\partial Q_{i}}{\partial p_{j}} = \frac{\partial^{2}\Phi}{\partial P_{j}\partial P_{i}}\text{ pour }i\le s\text{ et }j > s \\
                \frac{\partial P_{i}}{\partial p_{j}} = \frac{\partial^{2}\Phi}{\partial P_{j}\partial q_{i}}\text{ pour }i > s\text{ et }j > s
        \end{cases}
\eenn
Puisque nous avons des termes identiques car il y a juste une inversion des identifiants de lignes et de colonnes alors les d\'{e}terminants des deux matrices au num\'{e}rateur et au d\'{e}nominateur sont \'{e}gaux et nous avons donc bien $\lvert J\rvert = 1$.

Par extension, nous avons :
\be
	\int\mathrm{d}\Gamma = cste/\mathrm{t} \Leftrightarrow \frac{\mathrm{d}\Gamma}{\mathrm{dt}} = 0 \label{EQ:46_5}
\ee
L'invariance d'un volume de l'espace des phases par une transformation canonique est le \emph{th\'{e}or\`{e}me de Liouville}. Cela s'applique lors du mouvement puisque la variation au cours du temps de $p$ et $q$ est une transformation canonique de fonction g\'{e}n\'{e}ratrice $-S$ comme d\'{e}j\`{a} d\'{e}montr\'{e}e au paragraphe (\ref{PAR:45})

\section{\'{E}quations de Hamilton-Jacobi}\label{PAR:47}

\subsection{G\'{e}n\'{e}ralit\'{e}s}

Le paragraphe (\ref{PAR:43}) nous a permis de d\'{e}finir l'action comme fonction des coordonn\'{e}es et du temps. Nous avons aussi reli\'{e} la d\'{e}riv\'{e}e partielle par rapport au temps de $S(q,\mathrm{t})$ \`{a} la fonction de Hamilton du syst\`{e}me telle, voir l'\'{e}quation (\ref{EQ:43_5}) :
\benn
	\frac{\partial S}{\partial\mathrm{t}} + H(q,p,\mathrm{t}) = 0
\eenn
Sachant que les impulsions sont le r\'{e}sultat de la d\'{e}riv\'{e}e partielle par rapport aux coordonn\'{e}es de l'action, voir les relations (\ref{EQ:43_3}), nous obtenons :
\be
	\frac{\partial S}{\partial\mathrm{t}} + H(q_{1},\ldots,q_{s},\frac{\partial S}{\partial q_{1}},\ldots,\frac{\partial S}{\partial q_{s}}) = 0 \label{EQ:47_1}
\ee
Cette \'{e}quation du premier ordre est appel\'{e}e \emph{\'{e}quation de Hamilton-Jacobi}. Toute \'{e}quation aux d\'{e}riv\'{e}es partielles du premier ordre poss\`{e}de une solution d\'{e}pendant d'une fonction arbitraire. Cette solution est l'int\'{e}grale g\'{e}n\'{e}rale de l'\'{e}quation. En m\'{e}canique, cette int\'{e}grale g\'{e}n\'{e}rale de l'\'{e}quation de Hamilton-Jacobi ne joue pas une r\^{o}le essentiel au contraire de l'\emph{int\'{e}grale compl\`{e}te} qui contient autant de constantes arbitraires ind\'{e}pendantes qu'il y a de variables ind\'{e}pendantes.

Dans l'\'{e}quation de Hamilton-Jacobi, ces variables ind\'{e}pendantes sont les coordonn\'{e}es et le temps. Par la suite, pour un syst\`{e}me de $s$ degr\'{e} de libert\'{e}, l'int\'{e}grale compl\`{e}te de cette \'{e}quation doit contenir $s+1$ constantes arbitraires.

Puisque l'action $S$ ne figure dans l'\'{e}quation (\ref{EQ:47_1}) que par ses d\'{e}riv\'{e}es partielles, l'une des constantes arbitraires appara\^{i}t dans l'int\'{e}grale compl\`{e}te comme une grandeur additive. L'int\'{e}grale compl\`{e}te de l'\'{e}quation de Hamilton-Jacobi prend donc la forme suivante :
\be
	S = f(q_{1},\ldots,q_{s},\alpha_{1},\ldots,\alpha_{s},\mathrm{t}) + A \label{EQ:47_2}
\ee
avec $i \in \{1,\ldots,s\}$, $\alpha_{i}$ et $A$ des constantes arbitraires.

\subsection{Int\'{e}grale compl\`{e}te et solution des \'{e}quations du mouvement}\label{PAR:47_1}

R\'{e}alisons une transformation canonique des grandeurs $q$ et $p$ vers des nouvelles variables en choisissant comme fonction g\'{e}n\'{e}ratrice la fonction $f(\mathrm{t},q,\alpha)$ et comme nouvelle impulsions $\alpha_{1},\ldots,\alpha_{s}$. D\'{e}signons par $\beta_{1},\ldots,\beta_{s}$ les nouvelles coordonn\'{e}es. Puisque la fonction g\'{e}n\'{e}ratrice d\'{e}pend des anciennes coordonn\'{e}es et des nouvelles impulsions, nous utilisons les formules (\ref{EQ:43_8}) :
\benn
	p_{i} = \frac{\partial f}{\partial q_{i}}\text{, }\beta_{i} = \frac{\partial f}{\partial\alpha{i}}\text{ et }H' = H + \frac{\partial f}{\partial\mathrm{t}}
\eenn
Puisque la fonction g\'{e}n\'{e}ratrice satisfait \`{a} l'\'{e}quation de Hamilton-Jacobi alors cette derni\`{e}re nous permet d'\'{e}crire :
\benn
	\frac{\partial f}{\partial\mathrm{t}} + H = 0 \Rightarrow H' = 0
\eenn
Or les \'{e}quations du mouvement pour les nouvelles variables (\ref{EQ:45_3}) sont :
\benn
	\begin{cases}
		\dot{\beta}_{i} = \frac{\mathrm{d}\beta}{\mathrm{dt}} = \frac{\partial H'}{\partial\alpha_{i}} = 0 \\
		\dot{\alpha}_{i} = \frac{\mathrm{d}\alpha}{\mathrm{dt}} = -\frac{\partial H'}{\partial\beta_{i}} = 0
	\end{cases}
\eenn
soit :
\be
	\forall i\in\{1,\ldots,s\}\text{, }\alpha_{i} = cste\text{ et }\beta_{i} = cste \label{EQ:47_3}
\ee
De plus, les $s$ \'{e}quations $\beta_{i} = \frac{\partial f}{\partial\alpha_{i}}$ permettent d'exprimer les $s$ coordonn\'{e}es $q$ en fonction des $2s$ constantes arbitraires $\alpha$ et $\beta$ et nous obtenons l'int\'{e}grale g\'{e}n\'{e}rale des \'{e}quations du mouvement.

\subsection{M\'{e}thodologie}\label{PAR:47_2}

La solution du probl\`{e}me du mouvement d'un syst\`{e}me m\'{e}canique par la m\'{e}thode de Hamilton-Jacobi se r\'{e}alise avec les \'{e}tapes suivantes.

En partant de la fonction de Hamilton, nous \'{e}crivons l'\'{e}quation de Hamilton-Jacobi pour chercher l'int\'{e}grale compl\`{e}te (\ref{EQ:47_2}). En la d\'{e}rivant par rapport aux constantes arbitraires $\alpha$ et en \'{e}galant ces d\'{e}riv\'{e}es aux nouvelles constantes $\beta$, nous obtenons $s$ \'{e}quations alg\'{e}briques :
\be
	\frac{\partial S}{\partial\alpha_{i}} = \beta_{i} \label{EQ:47_4}
\ee
La solution de ce syst\`{e}me nous donne les coordonn\'{e}es $q$ en fonction du temps et des $2s$ constantes arbitraires. Les relations $\frac{\partial S}{\partial q_{i}} = p_{i}$ nous donnent ensuite la relation entre l'impulsion et le temps.

Dans le cas o\`{u} nous avons un syst\`{e}me conservatif, i.e. tel que la fonction de Hamilton ne d\'{e}pend explicitement pas du temps, alors la d\'{e}pendance de l'action par rapport au temps se r\'{e}duit au terme $-E\mathrm{t}$, voir la formule (\ref{EQ:44_2}), telle que :
\be
	S = S_{0}(q) - E\mathrm{t} \label{EQ:47_5}
\ee
En substituant cela dans l'\'{e}quation de Hamilton-Jacobi et puisque $\frac{\partial S}{\partial q_{i}} = \frac{\partial S_{0}}{\partial q_{i}}$, nous obtenons pour l'action r\'{e}duite :
\be
	H\left(q_{1},\ldots,q_{s},\frac{\partial S_{0}}{\partial q_{1}},\ldots,\frac{\partial S_{0}}{\partial q_{s}}\right) = E \label{EQ:47_6}
\ee

\section{S\'{e}paration des variables}\label{PAR:48}

\subsection{G\'{e}n\'{e}ralit\'{e}s}\label{PAR:48_1}

La m\'{e}thode de \emph{s\'{e}paration des variables} permet dans certains cas de trouver l'int\'{e}grale compl\`{e}te de l'\'{e}quation de Hamilton-Jacobi.

Nous admettons qu'une coordonn\'{e}e $q_{1}$ et que la d\'{e}riv\'{e}e $\frac{\partial S}{\partial q_{1}}$ n'entrent dans l'\'{e}quation de Hamilton-Jacobi que sous la forme d'une combinaison $\varphi\left(q_{1},\frac{\partial S}{\partial q_{1}}\right)$ ne contenant pas d'autres coordonn\'{e}es, ni d'autres d\'{e}riv\'{e}es, ni le temps. Autrement dit, nous admettons cons\'{e}cutivement que l'\'{e}quation soit sous la forme :
\be
	\forall i \in \{2,\ldots,s\}\text{, }\Phi\left\{q_{i},\mathrm{t},\frac{\partial S}{\partial q_{i}},\frac{\partial S}{\partial\mathrm{t}},\varphi\left(q_{1},\frac{\partial S}{\partial q_{1}}\right)\right\} = 0 \label{EQ:48_1}
\ee
Nous chercherons une solution sous la forme :
\be
	S = S'(q_{1},\ldots,q_{s},\mathrm{t}) + S_{1}(q_{1}) \label{EQ:48_2}
\ee
Comme :
\benn
	\frac{\partial S}{\partial q_{i}} = \frac{\partial S'}{\partial q_{i}}\text{, }\frac{\partial S}{\partial\mathrm{t}} = \frac{\partial S'}{\partial\mathrm{t}}\text{ et }\frac{\partial S}{\partial q_{1}} = \frac{\partial S_{1}}{\partial q_{1}} = \frac{\mathrm{d}S_{1}}{\mathrm{d}q_{1}}
\eenn
alors l'\'{e}quation (\ref{EQ:48_1}) devient :
\be
	\forall i \in \{2,\ldots,s\}\text{, }\Phi\left\{q_{i},\mathrm{t},\frac{\partial S'}{\partial q_{i}},\frac{\partial S'}{\partial\mathrm{t}},\varphi\left(q_{1},\frac{\mathrm{d}S_{1}}{\mathrm{d}q_{1}}\right)\right\} = 0 \label{EQ:48_3}
\ee
Comme l'action $S$ de (\ref{EQ:48_2}) est solution de (\ref{EQ:48_3}), l'identit\'{e} est valable quelque soit la valeur de $q_{1}$. Comme seule $\varphi$ d\'{e}pend de $q_{1}$ alors la relation (\ref{EQ:48_3}) n'est possible que si et seulement si $\varphi$ est constante. L'\'{e}quation (\ref{EQ:48_3}) peut alors se scinder en deux \'{e}quations :
\bea
	\varphi\left(q_{1},\frac{\mathrm{d}S_{1}}{\mathrm{d}q_{1}}\right) & = & \alpha_{1} \label{EQ:48_4} \\
	\Phi\left\{q_{i},\mathrm{t},\frac{\partial S'}{\partial q_{i}},\frac{\partial S'}{\partial\mathrm{t}},\alpha_{1}\right\} = 0 \label{EQ:48_5}
\eea
avec $\alpha_{1}$ une constante arbitraire.

L'\'{e}quation (\ref{EQ:48_4}) est une \'{e}quation diff\'{e}rentielle ordinaire \`{a} partir de laquelle la fonction $S_{1}(q_{1})$ peut \^{e}tre obtenue par int\'{e}gration. L'\'{e}quation (\ref{EQ:48_5}) reste une \'{e}quation aux d\'{e}riv\'{e}es partielles mais avec moins de variables ind\'{e}pendantes. S'il est possible de r\'{e}aliser la m\^{e}me s\'{e}paration successivement pour les $s$ coordonn\'{e}es et le temps, alors le calcul de l'int\'{e}grale compl\`{e}te de l'\'{e}quation de Hamilton-Jacobi se ram\`{e}nera \`{a} des calculs de primitives et d'int\'{e}grales.

Pour un syst\`{e}me conservatif, il s'agit de s\'{e}parer les $s$ coordonn\'{e}es de l'\'{e}quation (\ref{EQ:47_6}). L'int\'{e}grale cherch\'{e}e a la forme suivante :
\be
	S = \sum_{i}S_{i}(q_{i},\alpha_{1},\ldots,\alpha_{s}) - E(\alpha_{1},\ldots,\alpha_{s})\mathrm{t} \label{EQ:48_6}
\ee
o\`{u} la fonction $S_{i}$ ne d\'{e}pend que de la coordonn\'{e}e $q_{i}$ et o\`{u} l'\'{e}nergie $E$ est une fonction des constantes arbitraires $\alpha_{i}$. $E$ s'obtient en portant $S_{0} = \sum_{i}S_{i}$ dans l'\'{e}quation (\ref{EQ:48_6}).

\subsection{Cas particulier avec une variable cyclique}\label{PAR:48_2}

Nous prenons l'hypoth\`{e}se que la coordonn\'{e}e $q_{1}$ est cyclique, i.e. $\frac{\partial L}{\partial q_{1}} = 0$. La d\'{e}finition de la fonction de Hamilton (\ref{EQ:40_2}), $H = \sum_{i}p_{i}\dot{q}_{i} - L$ donne :
\benn
	\frac{\partial H}{\partial q_{1}} = 0 - \frac{\partial L}{\partial q_{1}} = 0
\eenn
Donc $q_{1}$ n'est pas une variable explicite de la fonction de Hamilton et elle n'entre donc pas dans l'\'{e}quation de Hamilton-Jacobi. De fait, la fonction $\varphi\left(q_{1},\frac{\partial S}{\partial q_{1}}\right)$ ne peut pas d\'{e}pendre de $q_{1}$. La fonction $\varphi$ peut alors s'\'{e}crire simplement $\frac{\partial S}{\partial q_{1}} = \frac{\mathrm{d}S_{1}}{\mathrm{d}q_{1}}$.

L'\'{e}quation (\ref{EQ:48_4}) donne alors :
\benn
	\frac{\mathrm{d}S_{1}}{\mathrm{d}q_{1}} = \alpha_{1} \Leftrightarrow S_{1} = \alpha_{1}q_{1}
\eenn
de telle sorte que la relation (\ref{EQ:48_2}) s'\'{e}crit alors :
\be
	S = S'(q_{i},\mathrm{t}) + \alpha_{1}q_{1}\text{ pour }i \ne 1 \label{EQ:48_7}
\ee
La relation (\ref{EQ:43_3}) nous donne $\frac{\partial S}{\partial q_{1}} = \frac{\partial S_{1}}{\partial q_{1}} = p_{1}$ ainsi $\alpha_{1}$ est ici la valeur constante de l'impulsion correspondant \`{a} la coordonn\'{e}e cyclique.

Dans un syst\`{e}me conservatif, nous avons par d\'{e}finition $\frac{\partial H}{\partial\mathrm{t}} = 0$. Alors la s\'{e}paration du temps dans le terme $-E\mathrm{t}$ correspond \`{a} la m\'{e}thode de s\'{e}paration des variables en prenant $\mathrm{t}$ pour variable cyclique.

\subsection{Exemples concrets}\label{PAR:48_3}

Pour s\'{e}parer les variables dans l'\'{e}quation de Hamilton-Jacobi, il est essentiel de bien choisir les coordonn\'{e}es. Nous \'{e}tudions par la suite plusieurs cas d'un point mat\'{e}riel de masse $m$ dans un syst\`{e}me conservatif.

\subsubsection{Coordonn\'{e}es sph\'{e}riques}\label{PAR:48_3_1}

L'exercice (\ref{PAR:40_EX3}) nous donne la fonction de Hamilton en fonction de $(r,\theta,\varphi)$ telle que :
\benn
	H = \frac{1}{2m}\left(p_{r}^{2} + \frac{p_{\theta}^{2}}{r^{2}} + \frac{p_{\varphi}^{2}}{r^{2}\sin^{2}\theta}\right) + U(r,\theta,\varphi)
\eenn
Les variables $r$, $\theta$ et $\varphi$ peuvent \^{e}tre s\'{e}par\'{e}es si et seulement si l'\'{e}nergie potentielle s'\'{e}crit sous la forme :
\benn
	U(r,\theta,\varphi) = a(r) + \frac{b(\theta)}{r^{2}} + \frac{c(\varphi)}{r^{2}\sin^{2}\theta}
\eenn
avec $a(r)$, $b(\theta)$ et $c(\varphi)$ des fonctions arbitraires.

La plupart des ph\'{e}nom\`{e}nes physiques ont une sym\'{e}trie axiale et dans ce cadre, nous pouvons donc r\'{e}duire la formulation de l'\'{e}nergie potentielle \`{a} : $U(r,\theta,\varphi) = a(r) + \frac{b(\theta)}{r^{2}}$. En utilisant les relations (\ref{EQ:43_3}), nous pouvons en d\'{e}duire la fonction de Hamilton :
\benn
	H = \frac{1}{2m}\left[\left(\frac{\partial S}{\partial r}\right)^{2} + \frac{1}{r^{2}}\left(\frac{\partial S}{\partial\theta}\right)^{2} + \frac{1}{r^{2}\sin^{2}\theta}\left(\frac{\partial S}{\partial\varphi}\right)^{2}\right] + a(r) + \frac{b(\theta)}{r^{2}}
\eenn
Donc l'\'{e}quation de Hamilton-Jacobi devient en utilisant la relation (\ref{EQ:47_5}) : $S = S_{0}(r,\theta,\varphi) - E\mathrm{t}$, soit $\frac{\partial S}{\partial\mathrm{t}} = E$ :
\bea
	& & \frac{\partial S}{\partial\mathrm{t}} + H = 0 \nonumber \\
	& \Leftrightarrow & \frac{1}{2m}\left(\frac{\partial S_{0}}{\partial r}\right)^{2} + a(r) + \frac{1}{2mr^{2}}\left[\left(\frac{\partial S_{0}}{\partial\theta}\right)^{2} + 2mb(\theta)\right] + \frac{1}{2mr^{2}\sin^{2}\theta}\left(\frac{\partial S_{0}}{\partial\varphi}\right)^{2} = E \nonumber
\eea
Dans notre hypoth\`{e}se de sym\'{e}trie axiale, la variable $\varphi$ est cyclique. Nous cherchons alors une solution telle que d\'{e}j\`{a} d\'{e}finie en (\ref{EQ:48_7}) : $S_{0} = p_{\varphi}\varphi + S_{1}(r) + S_{2}(\theta)$. Nous avons toujours :
\benn
	\frac{\mathrm{d}S_{1}}{\mathrm{d}r} = \frac{\partial S}{\partial r}\text{, }\frac{\mathrm{d}S_{1}}{\mathrm{d}\theta} = \frac{\partial S}{\partial\theta}\text{ et }\frac{\partial S}{\partial\varphi} = p_{\varphi}
\eenn
donc l'\'{e}quation de Hamilton-Jacobi devient :
\benn
	\frac{1}{2m}\left(\frac{\mathrm{d}S_{1}}{\mathrm{d}r}\right)^{2} + a(r) + \frac{1}{2mr^{2}}\left[\left(\frac{\mathrm{d}S_{2}}{\mathrm{d}\theta}\right)^{2} + 2mb(\theta)\right] + \frac{1}{2mr^{2}\sin^{2}\theta}p_{\varphi}^{2} = E
\eenn
Nous devons alors s\'{e}quentiellement r\'{e}soudre l'\'{e}quation :
\benn
	\left(\frac{\mathrm{d}S_{2}}{\mathrm{d}\theta}\right)^{2} + 2mb(\theta) + \frac{p_{\varphi}^{2}}{\sin^{2}\theta} = \beta
\eenn
puis :
\benn
	\frac{1}{2m}\left(\frac{\mathrm{d}S_{1}}{\mathrm{d}r}\right)^{2} + a(r) + \frac{\beta}{2mr^{2}} = E
\eenn
La premi\`{e}re donne :
\benn
	S_{2} = \int{\sqrt{\beta - 2mb(\theta) - \frac{p_{\varphi}^{2}}{\sin^{2}(\theta)}}\mathrm{d}\theta}
\eenn
et la seconde :
\benn
	S_{1} = \int{\sqrt{2m(E- a(r)) - \frac{\beta}{r^{2}}}\mathrm{d}r}
\eenn
ce qui nous permet d'obtenir pour l'action :
\be
	S = p_{\varphi}\varphi + \int{\sqrt{2m(E- a(r)) - \frac{\beta}{r^{2}}}\mathrm{d}r} + \int{\sqrt{\beta - 2mb(\theta) - \frac{p_{\varphi}^{2}}{\sin^{2}(\theta)}}\mathrm{d}\theta} - E\mathrm{t} \label{EQ:48_9}
\ee
qui est l'int\'{e}grale compl\`{e}te de l'\'{e}quation de Hamilton-Jacobi. Les constantes arbitraires sont $p_{\varphi}$, $\beta$ et $E$.

\subsubsection{Coordonn\'{e}es paraboliques}\label{PAR:48_3_2}

Le passage des coordonn\'{e}es cylindriques $(\rho,\varphi,z)$ aux coordonn\'{e}es paraboliques $(\xi,\eta,\varphi)$ s'op\`{e}re au moyen des formules de transformations suivantes :
\be
	z = \frac{1}{2}(\xi - \eta)\text{ et }\rho = \sqrt{\xi\eta} \label{EQ:48_10}
\ee
Les coordonn\'{e}es $\xi$ et $\eta$ appartiennent toutes les deux \`{a} $\mathbb{R}^{+}$. En utilisant le rayon $r$, nous avons :
\be
	r = \sqrt{z^{2} + \rho^{2}} = \sqrt{\frac{1}{4}(\xi^{2} + \eta^{2} - 2\xi\eta) + \xi\eta} = \frac{1}{2}(\xi + \eta) \label{EQ:48_11}
\ee
et encore :
\be
	\begin{cases}
		r + z = \frac{1}{2}(\xi + \eta) + \frac{1}{2}(\xi - \eta) = \xi \\
		r - z = \frac{1}{2}(\xi + \eta) - \frac{1}{2}(\xi - \eta) = \eta \label{EQ:48_12}
	\end{cases}
\ee
La formule (\ref{EQ:4_5}) donne la fonction de Lagrange en coordonn\'{e}es cylindriques \`{a} laquelle nous ajoutons la partie relative \`{a} l'\'{e}nergie potentielle :
\benn
	L = \frac{m}{2}(\dot{\rho}^{2} + \rho^{2}\dot{\varphi}^{2} + \dot{z}^{2}) - U(\rho,\varphi,z)
\eenn
En ins\'{e}rant la d\'{e}finition (\ref{EQ:48_10}) et parce que $\dot{z} = \frac{1}{2}(\dot{\xi} - \dot{\eta})$ et $\dot{\rho} = \frac{1}{2}\frac{\dot{\xi}\eta + \xi\dot{\eta}}{\sqrt{\xi\eta}}$, alors la fonction de Lagrange devient :
\bea
	L & = & \frac{m}{2}\left[\frac{1}{4}\frac{(\dot{\xi}\eta + \xi\dot{\eta})^{2}}{\xi\eta} + \xi\eta\dot{\varphi}^{2} + \frac{1}{4}(\xi^{2} - \eta^{2})^{2}\right] - U(\xi,\eta,\varphi) \nonumber \\
	& = & \frac{m}{8\xi\eta}(\dot{\xi}\eta^{2} + \xi^{2}\dot{\eta}^{2} + 2\xi\dot{\xi}\eta\dot{\eta} + \xi\dot{\xi}^{2}\eta + \xi\eta\dot{\eta}^{2} - 2\xi\dot{\xi}\eta\dot{\eta}) + \frac{m}{2}\xi\eta\dot{\varphi}^{2} - U(\xi,\eta,\varphi) \nonumber \\
	& = & \frac{m(\xi + \eta)}{8\xi\eta}(\dot{\xi}^{2}\eta + \xi\dot{\eta}^{2}) + \frac{m}{2}\xi\eta\dot{\varphi}^{2} - U(\xi,\eta,\varphi) \nonumber \\
	\Leftrightarrow L & = & \frac{m}{8}(\xi + \eta)\left(\frac{\dot{\xi}^{2}}{\xi} + \frac{\dot{\eta}^{2}}{\eta}\right) + \frac{m}{2}\xi\eta\dot{\varphi}^{2} - U(\xi,\eta,\varphi) \label{EQ:48_13}
\eea
En appliquant la d\'{e}finition de l'impulsion (\ref{EQ:7_5}), nous avons :
\benn
	p_{\xi} = \frac{\partial L}{\partial\dot{\xi}} = \frac{m}{4\xi}(\xi + \eta)\dot{\xi}\text{, }p_{\eta} = \frac{\partial L}{\partial\dot{\eta}} = \frac{m}{4\eta}(\xi + \eta)\dot{\eta}\text{ et }p_{\varphi} = m\xi\eta\dot{\varphi}
\eenn
La d\'{e}finition de la fonction de Hamilton (\ref{EQ:40_2}) nous donne ici :
\bea
	H & = & p_{\xi}\dot{\xi} + p_{\eta}\dot{\eta} + p_{\varphi}\dot{\varphi} - L \nonumber \\
	& = & \frac{4\xi p_{\xi}^{2}}{m(\xi + \eta)} + \frac{4\eta p_{\eta}^{2}}{m(\xi + \eta)} + \frac{p_{\varphi}^{2}}{m\xi\eta} \nonumber \\
	& & -\frac{m}{8}(\xi + \eta)\left[\frac{16\xi^{2}p_{\xi}^{2}}{m^{2}\xi(\xi + \eta)^{2}} + \frac{16\eta^{2}p_{\eta}^{2}}{m^{2}\eta(\xi + \eta)^{2}}\right] - \frac{m}{2}\xi\eta\frac{p_{\varphi}^{2}}{m^{2}\xi^{2}\eta^{2}} + U(\xi,\eta,\varphi) \nonumber \\
	& = & \frac{4\xi p_{\xi}^{2}}{m(\xi + \eta)} + \frac{4\eta p_{\eta}^{2}}{m(\xi + \eta)} - \frac{2\xi p_{\xi}^{2}}{m(\xi + \eta)} - \frac{2\eta p_{\eta}^{2}}{m(\xi + \eta)} + \frac{p_{\varphi}^{2}}{2m\xi\eta} + U(\xi,\eta,\varphi) \nonumber \\
	\Leftrightarrow H & = & \frac{2}{m}\frac{\xi p_{\xi}^{2} + \eta p_{\eta}^{2}}{(\xi + \eta)} \frac{p_{\varphi}^{2}}{2m\xi\eta} + U(\xi,\eta,\varphi) \label{EQ:48_14}
\eea
De la m\^{e}me mani\`{e}re que dans le paragraphe (\ref{PAR:48_3_1}), les syst\`{e}mes physiques ont g\'{e}n\'{e}ralement une sym\'{e}trie axiale. De plus, pour exercer la s\'{e}paration des variables, l'\'{e}nergie potentielle doit avoir une forme telle que :
\be
	U = \frac{a(\xi) + b(\eta)}{(\xi + \eta)} = \frac{a(r + z) + b(r - z)}{2r} \label{EQ:48_15}
\ee

Nous continuons d'utiliser la d\'{e}composition de l'action (\ref{EQ:47_5}) telle que $S = S_{0}(\xi,\eta,\varphi) - E\mathrm{t}$ qui nous donne $\frac{\partial S}{\partial\mathrm{t}} = -E$. Les relations (\ref{EQ:43_3}) nous permettent d'\'{e}crire :
\benn
	\frac{\partial S}{\partial\xi} = p_{\xi}\text{, }\frac{\partial S}{\partial\eta} = p_{\eta}\text{ et }\frac{\partial S}{\partial\varphi} = p_{\varphi}
\eenn
ou encore :
\benn
	\frac{\partial S_{0}}{\partial\xi} = p_{\xi}\text{, }\frac{\partial S_{0}}{\partial\eta} = p_{\eta}\text{ et }\frac{\partial S_{0}}{\partial\varphi} = p_{\varphi}
\eenn
L'\'{e}quation de Hamilton-Jacobi est alors dans notre cas :
\benn
	\frac{2}{m(\xi + \eta)}\left[\xi\left(\frac{\partial S_{0}}{\partial\xi}\right)^{2} + \eta\left(\frac{\partial S_{0}}{\partial\eta}\right)^{2}\right] + \frac{1}{2m\xi\eta}\left(\frac{\partial S_{0}}{\partial\varphi}\right)^{2} + \frac{a(\xi) + b(\eta)}{\xi + \eta} = E
\eenn
La coordonn\'{e}e $\varphi$ est cyclique par sym\'{e}trie axiale, elle se s\'{e}pare donc sous la forme $\frac{\partial S_{0}}{\partial\varphi}\varphi = p_{\varphi}\varphi$. L'\'{e}quation de Hamilton-Jacobi devient alors :
\bea
	2\xi\left(\frac{\partial S_{0}}{\partial\xi}\right)^{2} + 2\eta\left(\frac{\partial S_{0}}{\partial\eta}\right)^{2} + \frac{(\xi + \eta)}{2m\xi\eta}p_{\varphi}^{2} + ma(\xi) + mb(\eta)- mE(\xi + \eta) & = & 0 \nonumber \\
	\Leftrightarrow 2\xi\left(\frac{\partial S_{0}}{\partial\xi}\right)^{2} + ma(\xi) - mE\xi + \frac{p_{\varphi}^{2}}{2\xi} & & \nonumber \\
	+ 2\eta\left(\frac{\partial S_{0}}{\partial\eta}\right)^{2} + mb(\eta) - mE\eta + \frac{p_{\varphi}^{2}}{2\eta} & = & 0 \nonumber
\eea
En posant $S_{0} = p_{\varphi}\varphi + S_{1}(\xi) + S_{2}(\eta)$, cela implique $\frac{\partial S_{0}}{\partial\xi} = \frac{\partial S_{1}}{\partial\xi} = \frac{\mathrm{d}S_{1}}{\mathrm{d}\xi}$ et $\frac{\partial S_{0}}{\partial\eta} = \frac{\partial S_{2}}{\partial\xi} = \frac{\mathrm{d}S_{2}}{\mathrm{d}\xi}$. Nous avons donc \`{a} r\'{e}soudre s\'{e}quentiellement les deux \'{e}quations suivantes :
\bea
	2\xi\left(\frac{\mathrm{d}S_{1}}{\mathrm{d}\xi}\right)^{2} + ma(\xi) - mE\xi + \frac{p_{\varphi}^{2}}{2\xi} & = & \beta \nonumber \\
	2\eta\left(\frac{\mathrm{d}S_{2}}{\mathrm{d}\eta}\right)^{2} + mb(\eta) - mE\eta + \frac{p_{\varphi}^{2}}{2\eta} & = & -\beta \nonumber
\eea
Ce qui nous donne :
\bea
	S_{1} = \int{\sqrt{\frac{\beta}{2\xi} - \frac{p_{\varphi}^{2}}{4\xi^{2}} - \frac{m}{2\xi}a(\xi) + \frac{mE}{2}}\mathrm{d}\xi} \nonumber \\
	S_{2} = \int{\sqrt{-\frac{\beta}{2\eta} - \frac{p_{\varphi}^{2}}{4\eta^{2}} - \frac{m}{2\eta}a(\eta) + \frac{mE}{2}}\mathrm{d}\xi} \nonumber
\eea
L'action du syst\`{e}me en coordonn\'{e}es paraboliques, int\'{e}grale compl\`{e}te de l'\'{e}quation de Hamilton-Jacobi est :
\bea
	S & = & -E\mathrm{t} + p_{\varphi}\varphi + \int{\sqrt{\frac{\beta}{2\xi} - \frac{p_{\varphi}^{2}}{4\xi^{2}} - \frac{m}{2\xi}a(\xi) + \frac{mE}{2}}\mathrm{d}\xi} \nonumber \\
	& & + \int{\sqrt{-\frac{\beta}{2\eta} - \frac{p_{\varphi}^{2}}{4\eta^{2}} - \frac{m}{2\eta}a(\eta) + \frac{mE}{2}}\mathrm{d}\xi} \label{EQ:48_16}
\eea
avec $p_{\varphi}$, $\beta$ et $E$ des constantes arbitraires.

\subsubsection{Coordonn\'{e}es elliptiques}\label{PAR:48_3_3}

Les coordonn\'{e}es elliptiques $(\xi,\eta,\varphi)$ d\'{e}coulent des formules :
\be
	\begin{cases}
		\rho = \sigma\sqrt{(\xi^{2} - 1)(1 - \eta^{2})} \nonumber \\
		z = \sigma\xi\eta \label{EQ:48_17}
	\end{cases}
\ee
La constante $\sigma$ est le param\`{e}tre de la transformation. $\xi$ varie dans $[1;+\infty[$ et $\eta$ dans $[-1;1]$. D'un point de vue g\'{e}om\'{e}trique, la quantit\'{e} $r_{1} = \sqrt{(z - \sigma)^{2} + \rho^{2}}$, respectivement $r_{2} = \sqrt{(z - \sigma)^{2} + \rho^{2}}$, est la distance au point $A_{1}$, respectivement $A_{2}$, de l'axe $z$, qui a pour coordonn\'{e}e $z = \sigma$, respectivement $z = -\sigma$. Nous avons alors :
\be
	\begin{cases}
		r_{1} = \sqrt{(\sigma\xi\eta - \sigma)^{2} + \sigma^{2}(\xi^{2} - 1)(1 - \eta^{2})} = \sigma\sqrt{\xi^{2}\eta^{2} + 1 - 2\xi\eta + \xi^{2} - \xi^{2}\eta^{2} - 1 - \eta^{2}} \nonumber \\
		= \sigma\sqrt{\xi^{2} + \eta^{2} - 2\xi\eta} = \sigma(\xi - \eta) \nonumber \\
		r_{1} = \sqrt{(\sigma\xi\eta + \sigma)^{2} + \sigma^{2}(\xi^{2} - 1)(1 - \eta^{2})} = \sigma\sqrt{\xi^{2}\eta^{2} + 1 + 2\xi\eta + \xi^{2} - \xi^{2}\eta^{2} - 1 - \eta^{2}} \nonumber \\
		= \sigma\sqrt{\xi^{2} + \eta^{2} + 2\xi\eta} = \sigma(\xi + \eta) \nonumber \\
		\xi = \frac{r_{2} + r_{1}}{2\sigma}\text{ et }\eta = \frac{r_{2} - r_{1}}{2\sigma} \label{EQ:48_18}
	\end{cases}
\ee
La fonction de Lagrange en coordonn\'{e}es cylindriques est, voir la formule (\ref{EQ:4_5}) :
\benn
	L = \frac{m}{2}(\dot{r}^{2} + r^{2}\dot{\varphi}^{2} + \dot{z}^{2}) - R(r,\varphi,z)
\eenn
o\`{u} la coordonn\'{e}e $r$ est identifiable \`{a} $\rho$, alors :
\bea
	\dot{r} & = & \dot{\rho} = \sigma\frac{\left(2\xi\dot{\xi}(1 - \eta^{2}) - 2\eta\dot{\eta}(\xi^{2} - 1)\right)}{2\sqrt{(\xi^{2} - 1)(1 - \eta^{2})}} = \sigma\frac{\left(\xi\dot{\xi}(1 - \eta^{2}) - \eta\dot{\eta}(\xi^{2} - 1)\right)}{\sqrt{(\xi^{2} - 1)(1 - \eta^{2})}} \nonumber \\
	\dot{z} & = & \sigma(\dot{\xi}\eta + \xi\dot{\eta}) \nonumber
\eea
et la fonction de Lagrange se d\'{e}veloppe :
\benn
	L = \frac{m}{2}\left[\sigma^{2}\frac{\left(\xi\dot{\xi}(1 - \eta^{2}) - \eta\dot{\eta}(\xi^{2} - 1)\right)^{2}}{(\xi^{2} - 1)(1 - \eta^{2})} + \sigma^{2}(\xi^{2} - 1)(1 - \eta^{2})\dot{\varphi}^{2} + \sigma^{2}(\dot{\xi}\eta + \xi\dot{\eta})^{2}\right] - U(\xi,\eta\varphi)
\eenn
Il nous est n\'{e}cessaire de faire un focus sur la quantit\'{e} :
\bea
	\left(\xi\dot{\xi}(1 - \eta^{2} - \eta\dot{\eta}(\xi^{2} - 1)\right)^{2} & = & \xi^{2}\dot{\xi}^{2}(1 - \eta^{2}) + \eta^{2}\dot{\eta}^{2}(\xi^{2} - 1)^{2} - 2\xi\dot{\xi}\eta\dot{\eta}(\xi^{2} - 1)(1 - \eta^{2}) \nonumber \\
	& & + (\dot{\xi}^{2}\eta^{2} + \xi^{2}\dot{\eta}^{2} + 2\xi\dot{\xi}\eta\dot{\eta})(\xi^{2} - 1)(1 - \eta^{2}) \nonumber \\
	& = & \xi^{2}\dot{\xi}^{2} + \xi^{2}\dot{\xi}^{2}\eta^{4} - 2\xi^{2}\dot{\xi}^{2}\eta^{2} + \eta^{2}\dot{\eta}^{2}\xi^{4} + \eta^{2}\dot{\eta}^{2} - 2\xi^{2}\eta^{2}\dot{\eta}^{2} \nonumber \\
	& & + (\dot{\xi}{2}\eta^{2} + \xi^{2}\dot{\eta}^{2})(\xi^{2} + \eta^{2} - \xi^{2}\eta^{2} - 1) \nonumber \\
	& = & \xi^{2}\dot{\xi}^{2} + \xi^{2}\dot{\xi}^{2}\eta^{4} - 2\xi^{2}\dot{\xi}^{2}\eta^{2} + \eta^{2}\dot{\eta}^{2}\xi^{4} + \eta^{2}\dot{\eta}^{2} - 2\xi^{2}\eta^{2}\dot{\eta}^{2} \nonumber \\
	& & \xi^{2}\dot{\xi}^{2}\eta^{2} + \dot{\xi}^{2}\eta^{4} - \xi^{2}\dot{\xi}^{2}\eta^{4} - \dot{\xi}^{2}\eta^{2} + \xi^{4}\dot{\xi}^{2} + \xi^{2}\eta^{2}\dot{\xi}^{2} - \xi^{4}\eta^{2}\dot{\eta}^{2} - \xi^{2}\dot{\eta}^{2} \nonumber \\
	& = & \dot{\xi}^{2}(\xi^{2} - 2\xi^{2}\eta^{2} + \xi^{2}\eta^{2} + \eta^{4} - \eta^{2}) + \dot{\eta}^{2}(\eta^{2} - 2\xi^{2}\eta^{2} + \xi^{4} + \xi^{2}\eta^{2} - \xi^{2}) \nonumber \\
	& = & \dot{\xi}^{2}(\eta^{4} - \eta^{2} - \xi^{2}\eta^{2} + \xi^{2}) + \dot{\eta}^{2}(\eta^{4} - \eta^{2} - \xi^{2}\eta^{2} + \eta^{2}) \nonumber \\
	& = & \dot{\xi}^{2}(1 - \eta^{2})(\xi^{2} - \eta^{2}) + \dot{\eta}^{2}(\xi^{2} - 1)(\xi^{2} - \eta^{2}) \nonumber
\eea
qui permet d'\'{e}crire la fonction de Lagrange et de la d\'{e}velopper ainsi :
\bea
	L & = & \frac{m\sigma^{2}}{2}\left[\frac{\dot{\xi}^{2}(1 - \eta^{2})(\xi^{2} - \eta^{2}) + \dot{\eta}^{2}(\xi^{2} - 1)(\xi^{2} - \eta^{2})}{(\xi^{2} - 1)(1 - \eta^{2})}\right] + \frac{m\sigma^{2}}{2}(\xi^{2} - 1)(1 - \eta^{2})\dot{\varphi}^{2} - U(\xi,\eta,\varphi) \nonumber \\
	\Leftrightarrow L & = & \frac{m\sigma^{2}}{2}(\xi^{2} - \eta^{2})\left(\frac{\dot{\xi}^{2}}{(\xi^{2} - 1)} + \frac{\dot{\eta}^{2}}{(1 - \eta^{2})}\right) + \frac{m\sigma^{2}}{2}(\xi^{2} - 1)(1 - \eta^{2})\dot{\varphi}^{2} - U(\xi,\eta,\varphi) \label{EQ:48_19}
\eea
La d\'{e}finition de la fonction de Hamilton (\ref{EQ:40_2}) donne :
\bea
	H & = & p_{\xi}\dot{\xi} + p_{\eta}\dot{\eta} + p_{\varphi}\dot{\varphi} - L \nonumber \\
	& = & \frac{(\xi^{2} - 1)}{m\sigma^{2}(\xi^{2} - \eta^{2})}p_{\xi}^{2} + \frac{(1 - \eta^{2})}{m\sigma^{2}(\xi^{2} - \eta^{2})}p_{\eta}^{2} + \frac{p_{\varphi}^{2}}{m\sigma^{2}(\xi^{2} - 1)(1 - \eta^{2})} - \frac{m\sigma^{2}(\xi^{2} - \eta^{2})}{2(\xi^{2} - 1)}\frac{(\xi^{2} - 1)^{2}}{m^{2}\sigma^{4}(\xi^{2} - \eta^{2})^{2}}p_{\xi}^{2} \nonumber \\
	& & -\frac{m\sigma^{2}(\xi^{2} - \eta^{2})}{2(1 - \eta^{2})}\frac{(1 - \eta^{2})^{2}}{m^{2}\sigma^{4}(\xi^{2} - \eta^{2})^{2}}p_{\eta}^{2} - \frac{m\sigma^{2}}{2}(\xi^{2} - 1)(1 - \eta^{2})\frac{p_{\varphi}^{2}}{m^{2}\sigma^{4}(\xi^{2} - 1)^{2}(1 - \eta^{2})^{2}} + U(\xi,\eta,\varphi) \nonumber \\
	& = & \frac{(\xi^{2} - 1)}{2m\sigma^{2}(\xi^{2} - \eta^{2})}p_{\xi}^{2} + \frac{(1 - \eta^{2})}{2m\sigma^{2}(\xi^{2} - \eta^{2})}p_{\eta}^{2} + \frac{p_{\varphi}^{2}}{2m\sigma^{2}(\xi^{2} - 1)(1 - \eta^{2})} + U(\xi,\eta,\varphi) \nonumber
\eea
Or :
\benn
	\frac{1}{\xi^{2} - 1} + \frac{1}{1 - \eta^{2}} = \frac{\xi^{2} - \eta^{2}}{(\xi^{2} - 1)(1 - \eta^{2})}
\eenn
Donc la fonction de Hamilton devient :
\be
	H = \frac{1}{2m\sigma^{2}(\xi^{2} - \eta^{2})}\left[(\xi^{2} - 1)p_{\xi}^{2} + (1 - \eta^{2})p_{\eta}^{2} + \left(\frac{1}{\xi^{2} - 1} + \frac{1}{1 - \eta^{2}}\right)p_{\varphi}^{2}\right] + U(\xi,\eta,\varphi) \label{EQ:48_20}
\ee
En s'appuyant sur la sym\'{e}trie tr\`{e}s g\'{e}n\'{e}ralement axiale des ph\'{e}nom\`{e}nes physiques, i.e $\varphi$ est une variable cyclique, et en appliquant la m\'{e}thode de s\'{e}paration des variables, nous pouvons supposer que l'\'{e}nergie potentielle peut s'\'{e}crire sous la forme :
\be
	U = \frac{a(\xi) + b(\eta)}{\xi^{2} - \eta^{2}} = \frac{\sigma^{2}}{r_{1}r_{2}}\left(a\left(\frac{r_{2} + r_{1}}{2}\right) + b\left(\frac{r_{2} - r_{1}}{2}\right)\right) \label{EQ:48_21}
\ee
o\`{u} $a(\xi)$ et $b(\eta)$ sont des fonctions arbitraires.

En reprenant les m\^{e}mes hypoth\`{e}ses que dans le paragraphe (\ref{PAR:48_3_2}), \`{a} savoir $S = S_{0}(\xi,\eta,\varphi) - E\mathrm{t}$ et donc :
\benn
	p_{\xi} = \frac{\partial S}{\partial\xi} = \frac{\partial S_{0}}{\partial\xi}\text{, }p_{\eta} = \frac{\partial S}{\partial\eta} = \frac{\partial S_{0}}{\partial\eta}\text{, }p_{\varphi} = \frac{\partial S}{\partial\varphi} = \frac{\partial S_{0}}{\partial\varphi}\text{ et }\frac{\partial S}{\partial\mathrm{t}} = -E
\eenn
alors l'\'{e}quation de Hamilton-Jacobi s'\'{e}crit :
\bea
	\frac{1}{2m\sigma^{2}(\xi^{2} - \eta^{2})}\left[(\xi^{2} - 1)\left(\frac{\partial S_{0}}{\partial\xi}\right)^{2} + (1 - \eta^{2})\left(\frac{\partial S_{0}}{\partial\eta}\right)^{2} + \left(\frac{1}{\xi^{2} - 1} + \frac{1}{1 - \eta^{2}}\right)p_{\varphi}^{2}\right] + \frac{a(\xi) + b(\eta)}{\xi^{2} - \eta^{2}} & = & E \nonumber \\
	\Leftrightarrow (\xi^{2} - 1)\left(\frac{\partial S_{0}}{\partial\xi}\right)^{2} + \left(\frac{1}{\xi^{2} - 1}\right)p_{\varphi}^{2} + 2m\sigma^{2}a(\xi) - 2m\sigma^{2}\xi^{2}E & & \nonumber \\
	+(1 - \eta^{2})\left(\frac{\partial S_{0}}{\partial\eta}\right)^{2} + \left(\frac{1}{1 - \eta^{2}}\right)p_{\varphi}^{2} + 2m\sigma^{2}b(\eta) + 2m\sigma^{2}\eta^{2}E & = & 0 \nonumber
\eea
Nous posons alors $S_{0} = p_{\varphi}\varphi + S_{1}(\xi) + S_{2}(\eta)$ qui nous permet d'\'{e}crire :
\benn
	\frac{\partial S_{0}}{\partial\xi} = \frac{\partial S_{1}}{\partial\xi} = \frac{\mathrm{d}S_{1}}{\mathrm{d}\xi}\text{ et }\frac{\partial S_{0}}{\partial\eta} = \frac{\partial S_{2}}{\partial\xi} = \frac{\mathrm{d}S_{2}}{\mathrm{d}\xi}
\eenn
Nous avons donc \`{a} r\'{e}soudre s\'{e}quentiellement les deux \'{e}quations suivantes :
\bea
	(\xi^{2} - 1)\left(\frac{\mathrm{d}S_{1}}{\mathrm{d}\xi}\right)^{2} + \left(\frac{1}{\xi^{2} - 1}\right)p_{\varphi}^{2} + 2m\sigma^{2}a(\xi) - 2m\sigma^{2}\xi^{2}E & = & \beta \nonumber \\
	(1 - \eta^{2})\left(\frac{\mathrm{d}S_{2}}{\mathrm{d}\xi}\right)^{2} + \left(\frac{1}{1 - \eta^{2}}\right)p_{\varphi}^{2} + 2m\sigma^{2}b(\eta) + 2m\sigma^{2}\eta^{2}E & = & -\beta \nonumber
\eea
Nous obtenons donc :
\bea
	S_{1} & = & \int{\sqrt{\frac{2m\sigma^{2}\xi^{2}}{\xi^{2} - 1}E + \frac{\beta - 2m\sigma^{2}a(\xi)}{\xi^{2} - 1} - \frac{p_{\varphi}^{2}}{(\xi^{2} - 1)^{2}}}\mathrm{d}\xi} \nonumber \\
	S_{2} & = & \int{\sqrt{\frac{2m\sigma^{2}\eta^{2}}{1 - \eta^{2}}E - \frac{\beta + 2m\sigma^{2}b(\eta)}{1 - \eta^{2}} - \frac{p_{\varphi}^{2}}{(1 - \eta^{2})^{2}}}\mathrm{d}\eta} \nonumber
\eea
Le r\'{e}sultat de l'\'{e}quation de Hamilton-Jacobi est donc l'int\'{e}grale compl\`{e}te suivante\footnote{Nous obtenons ici une expression diff\'{e}rente par rapport au livre sur le facteur multiplicatif de l'\'{e}nergie.} :
\bea
	S & = & -E\mathrm{t} + p_{\varphi}\varphi + \int{\sqrt{\frac{2m\sigma^{2}\xi^{2}}{\xi^{2} - 1}E + \frac{\beta - 2m\sigma^{2}a(\xi)}{\xi^{2} - 1} - \frac{p_{\varphi}^{2}}{(\xi^{2} - 1)^{2}}}\mathrm{d}\xi} \nonumber \\
	& & + \int{\sqrt{\frac{2m\sigma^{2}\eta^{2}}{1 - \eta^{2}}E - \frac{\beta + 2m\sigma^{2}b(\eta)}{1 - \eta^{2}} - \frac{p_{\varphi}^{2}}{(1 - \eta^{2})^{2}}}\mathrm{d}\eta} \label{EQ:48_22}
\eea

\section{Invariant adiabatiques}\label{PAR:49}

\subsection{Cas g\'{e}n\'{e}ral}\label{PAR:49_1}

Nous consid\'{e}rons un syst\`{e}me m\'{e}canique anim\'{e} d'un mouvement lin\'{e}aire fini et caract\'{e}ris\'{e} par un param\`{e}tre $\lambda$ qui d\'{e}finit les propri\'{e}t\'{e}s du syst\`{e}me lui-m\^{e}me ou d'un champ ext\'{e}rieur dans lequel le syst\`{e}me se trouve.

Nous faisons l'hypoth\`{e}se que $\lambda$ est \emph{adiabatique}, i.e. qu'il varie peu au cours d'une p\'{e}riode $T$ du mouvement du syst\`{e}me lors d'une transformation. Hypoth\`{e}se que nous pouvons exprimer ainsi :
\be
	T\frac{\mathrm{d}\lambda}{\mathrm{dt}} \ll \lambda \label{EQ:49_1}
\ee
Avec $\lambda$ constant, le syst\`{e}me serait ferm\'{e} et effectuerait un mouvement strictement p\'{e}riodique d'\'{e}nergie constante $E$ et de p\'{e}riode $T(E)$ parfaitement d\'{e}termin\'{e} par le paragraphe (\ref{PAR:11}). Lorsque $\lambda$ est variable, le syst\`{e}me n'est pas ferm\'{e} et son \'{e}nergie ne se conserve pas.

Puisque $\lambda$ varie lentement ici, la vitesse de variation de l'\'{e}nergie $\dot{E}$ sera faible \'{e}galement. En prenant la moyenne de cette vitesse sur une p\'{e}riode $T$, les variations rapides en sont liss\'{e}es et la valeur $\overline{\dot{E}}$ d\'{e}termine alors la vitesse de variation lente de l'\'{e}nergie du syst\`{e}me. Nous pouvons en d\'{e}duire que $\overline{\dot{E}}$ sera proportionnelle \`{a} la vitesse $\dot{\lambda}$. Donc la grandeur $E$ varie dans le sens indiqu\'{e} par la variation de $\lambda$ et la d\'{e}pendance de $E$ par rapport \`{a} $\lambda$ peut \^{e}tre traduite comme une combinaison de $E$ et de $\lambda$. Cette combinaison qui reste constante pendant le mouvement d'un syst\`{e}me pour lequel les param\`{e}tres varient lentement se nomme \emph{invariant adiabatique}.

Soit $H(p,q,\lambda)$ la fonction de Hamilton du syst\`{e}me d\'{e}pendant du param\`{e}tre $\lambda$. L'\'{e}quation (\ref{EQ:40_5}) nous donne :
\bea
	\frac{\mathrm{d}H}{\mathrm{dt}} & = & \frac{\partial H}{\partial\mathrm{t}} \Leftrightarrow \frac{\mathrm{d}E}{\mathrm{dt}} = \frac{\partial H}{\partial\mathrm{t}} \nonumber \\
	\frac{\mathrm{d}E}{\mathrm{dt}} & = & \frac{\partial H}{\partial\lambda}\frac{\mathrm{d}\lambda}{\mathrm{dt}} \label{EQ:49_2}
\eea
car $\lambda$ ne d\'{e}pend pas explicitement du temps.

Mais la fonction de Hamilton d\'{e}pend des coordonn\'{e}es et des impulsions qui sont toutes les deux rapidement variables. Pour rendre plus \'{e}vidente la variation lente de l'\'{e}nergie, nous allons moyenner la relation (\ref{EQ:49_3}) sur une p\'{e}riode du mouvement telle que :
\benn
	\overline{\frac{\mathrm{d}E}{\mathrm{dt}}} = \overline{\left(\frac{\partial H}{\partial\lambda}\frac{\mathrm{d}\lambda}{\mathrm{dt}}\right)} = \overline{\frac{\partial H}{\partial\lambda}}\frac{\mathrm{d}\lambda}{\mathrm{dt}} + \frac{\partial H}{\partial\lambda}\overline{\frac{\mathrm{d}\lambda}{\mathrm{dt}}}
\eenn
Or, nous avons $\overline{\frac{\mathrm{d}\lambda}{\mathrm{dt}}}\approx 0$ car $\lambda$ varie lentement donc :
\be
	\overline{\frac{\mathrm{d}E}{\mathrm{dt}}} = \overline{\frac{\partial H}{\partial\lambda}}\frac{\mathrm{d}\lambda}{\mathrm{dt}} \label{EQ:49_3}
\ee
Encore une fois, $\lambda$ variant peu sur une p\'{e}riode du mouvement du syst\`{e}me, la quantit\'{e} $\overline{\frac{\partial H}{\partial\lambda}}$ peut \^{e}tre consid\'{e}r\'{e}e comme fonction de $p$ et $q$ uniquement. La valeur de cette quantit\'{e} se calcule ainsi :
\benn
	\overline{\frac{\partial H}{\partial\lambda}} = \frac{1}{T}\int_{0}^{T}{\frac{\partial H}{\partial\lambda}\mathrm{dt}}
\eenn
Or les \'{e}quations de Hamilton (\ref{EQ:40_4}) nous permettent d'\'{e}crire :
\benn
	\dot{q} = \frac{\mathrm{d}q}{\mathrm{dt}} = \frac{\partial H}{\partial p} \Leftrightarrow \mathrm{dt} = \left(\frac{\partial H}{\partial p}\right)^{-1}\mathrm{d}q
\eenn
qui, dans le calcul de la p\'{e}riode $T$, nous donne :
\be
	T = \int_{0}^{T}{\mathrm{dt}} = \oint{\left(\frac{\partial H}{\partial p}\right)^{-1}\mathrm{d}q} \label{EQ:49_4}
\ee
o\`{u} le symbole $\oint$ d\'{e}signe l'int\'{e}gration par rapport \`{a} la variation totale de la coordonn\'{e}e au cours de la p\'{e}riode. Par exemple, si le syst\`{e}me effectue une rotation et si $q$ est l'angle du mouvement de rotation, alors l'int\'{e}gration par rapport \`{a} $\mathrm{d}q$ est r\'{e}alis\'{e}e entre $0$ et $2\pi$.

La relation (\ref{EQ:49_3}) peut se r\'{e}\'{e}crire ainsi :
\be
	\overline{\frac{\mathrm{d}E}{\mathrm{dt}}} = \frac{\mathrm{d}\lambda}{\mathrm{dt}}\frac{\oint\left(\frac{\partial H}{\partial\lambda}\right)\left(\frac{\partial H}{\partial p}\right)^{-1}\mathrm{d}q}{\oint{\left(\frac{\partial H}{\partial p}\right)^{-1}\mathrm{d}q}} \label{EQ:49_5}
\ee
Comme indiqu\'{e} pr\'{e}c\'{e}demment, les int\'{e}grales sont prises le long de la trajectoire, $\lambda$ \'{e}tant consid\'{e}r\'{e} comme constant. Le long de la trajectoire, la fonction de Hamilton conserve l'\'{e}nergie $E$ constante, puisque le param\`{e}tre $\lambda$ est consid\'{e}r\'{e} comme constant, et l'impulson est une fonction d\'{e}finie de la coordonn\'{e}e variable $q$ et des deux param\`{e}tres $\lambda$ et $E$.

La conservation de l'\'{e}nergie nous donne :
\bea
	H(p,q,\lambda) & = & E \Leftrightarrow \frac{\mathrm{d}H(h,p,\lambda)}{\mathrm{d}\lambda} = 0 \Leftrightarrow \frac{\partial H}{\partial\lambda} + \frac{\partial H}{\partial p}\frac{\partial p}{\partial\lambda} = 0 \nonumber \\
	\Leftrightarrow \frac{\partial p}{\partial\lambda} & = & -\frac{\partial H}{\partial\lambda}\left(\frac{\partial H}{\partial p}\right)^{-1}
\eea
Nous portons ce r\'{e}sultat dans la relation (\ref{EQ:49_5}) pour obtenir :
\benn
	\overline{\frac{\mathrm{d}E}{\mathrm{dt}}} = -\frac{\mathrm{d}\lambda}{\mathrm{dt}}\frac{\oint{\frac{\partial p}{\partial\lambda}\mathrm{d}q}}{\oint{\frac{\partial p}{\partial E}\mathrm{d}q}}
\eenn
que nous pouvons d\'{e}velopper ainsi :
\benn
	\overline{\frac{\mathrm{d}E}{\mathrm{dt}}}\oint{\frac{\partial p}{\partial E}\mathrm{d}q} + \frac{\mathrm{d}\lambda}{\mathrm{dt}}\oint{\frac{\partial p}{\partial\lambda}\mathrm{d}q} \Leftrightarrow \oint{\left(\overline{\frac{\mathrm{d}E}{\mathrm{dt}}}\frac{\partial p}{\partial E} + \frac{\mathrm{d}\lambda}{\mathrm{dt}}\frac{\partial p}{\partial\lambda}\right)\mathrm{d}q} = 0
\eenn
Cette derni\`{e}re \'{e}galit\'{e} peut aussi s'\'{e}crire :
\be
	\overline{\frac{\mathrm{d}I}{\mathrm{dt}}} = 0 \label{EQ:49_6}
\ee
De plus comme l'impulsion est fonction de $q$, $\lambda$ et $E$ alors :
\benn
	\dot{p} = \frac{\mathrm{d}p}{\mathrm{dt}} = \frac{\partial p}{\partial E}\frac{\mathrm{d}E}{\mathrm{dt}} + \frac{\partial p}{\partial\lambda}\frac{\mathrm{d}\lambda}{\mathrm{dt}} + \frac{\partial p}{\partial q}\frac{\mathrm{d}q}{\mathrm{dt}} \nonumber
\eenn
En utilisant les \'{e}quations de Hamilton (\ref{EQ:40_4}), $\dot{q} = \frac{\partial H}{\partial p}$ et $\dot{p} = -\frac{\partial H}{\partial q}$, le relation se d\'{e}veloppe ainsi :
\bea
	\dot{p} & = & \frac{\partial p}{\partial E}\frac{\mathrm{d}E}{\mathrm{dt}} + \frac{\partial p}{\partial\lambda}\frac{\mathrm{d}\lambda}{\mathrm{dt}} + \frac{\partial p}{\partial q}\frac{\partial H}{\partial p} = \frac{\partial p}{\partial E}\frac{\mathrm{d}E}{\mathrm{dt}} + \frac{\partial p}{\partial\lambda}\frac{\mathrm{d}\lambda}{\mathrm{dt}} - \dot{p} \nonumber \\
	\Leftrightarrow 2\dot{p} & = & \frac{\partial p}{\partial E}\frac{\mathrm{d}E}{\mathrm{dt}} + \frac{\partial p}{\partial\lambda}\frac{\mathrm{d}\lambda}{\mathrm{dt}} \nonumber \\
	\Leftrightarrow \overline{\frac{\mathrm{d}p}{\mathrm{dt}}} & = & \frac{\partial p}{\partial E}\overline{\frac{\mathrm{d}E}{\mathrm{dt}}} + \frac{\partial p}{\partial\lambda}\frac{\mathrm{d}\lambda}{\mathrm{dt}} \nonumber
\eea
Cette derni\`{e}re relation permet d'identifier et de d\'{e}finir la quantit\'{e} :
\be
	I = \frac{1}{2\pi}\oint{p\mathrm{d}q} \label{EQ:49_7}
\ee
en ajoutant un facteur de normalisation $\pi$. $I$ est une int\'{e}grale prise le long de la trajectoire pour $E$ et $\lambda$ donn\'{e}es.

Dans l'approximation consid\'{e}r\'{e}e, la quantit\'{e} $I$ reste donc constante quand $\lambda$ varie. $I$ est donc un avariant adiabatique qui est fonction de l'\'{e}nergie du syst\`{e}me et du param\`{e}tre $\lambda$.

\`{A} partir de ces nouveaux \'{e}l\'{e}ments, la p\'{e}riode $T$, voir la formule (\ref{EQ:49_4}) peut \^{e}tre \'{e}crite diff\'{e}remment :
\be
	T = \oint{\left(\frac{\partial H}{\partial p}\right)^{-1}\mathrm{d}q} = \oint{\left(\frac{\partial E}{\partial p}\right)\mathrm{d}q} = \oint{\frac{\partial p}{\partial E}\mathrm{d}q} = -2\pi\frac{\partial I}{\partial E} \label{EQ:49_8}
\ee
Or par d\'{e}finition, la pulsation se d\'{e}finit comme $\omega = 2\pi/T$ donc :
\be
	\omega = \left(\frac{\partial I}{\partial E}\right)^{-1} = \frac{\partial E}{\partial I} \label{EQ:49_9}
\ee
Ainsi la d\'{e}riv\'{e}e de l'invariant adiabatique $I$ par rapport \`{a} l'\'{e}nergie d\'{e}termine la p\'{e}riode du mouvement.

\subsection{Cas \`{a} un unique degr\'{e} de libert\'{e}}\label{PAR:49_2}

L'int\'{e}grale (\ref{EQ:49_7}) a une interpr\'{e}tation g\'{e}om\'{e}trique dans l'espace des phases. Dans le cas \`{a} un seul degr\'{e} de libert\'{e}, l'espace des phase est le syst\`{e}me de coordonn\'{e}es $(p,q)$ \`{a} deux dimensions. Pour un mouvement p\'{e}riodique, la trajectoire dans l'espace des phases est une courbe ferm\'{e}e. L'int\'{e}grale dans l'\'{e}galit\'{e} (\ref{EQ:49_7}) prise le long de cette courbe repr\'{e}sente la surface int\'{e}rieure et l'invariant adiabatique se calcule aussi de cette mani\`{e}re :
\be
	I = \frac{1}{2\pi}\iint{\mathrm{d}p\mathrm{d}q} \label{EQ:49_10}
\ee
Nous prenons \`{a} partir de maintenant l'exemple de la d\'{e}termination de l'invariant adiabatique pour un oscillateur lin\'{e}aire. Ce dernier a \'{e}t\'{e} \'{e}tudi\'{e} au paragraphe (\ref{PAR:21}) et l'\'{e}nergie calcul\'{e}e nous permet d'\'{e}crire la fonction de Hamilton dans le cas de la conservation d'\'{e}nergie :
\be
	H = \frac{p^{2}}{2m} + \frac{m\omega^{2}}{2}q^{2} \label{EQ:49_11}
\ee
avec $\omega$ la fr\'{e}quence propre du syst\`{e}me. La trajectoire dans l'espace des phases est donn\'{e}e par :
\benn
	H(p,q) = E \Leftrightarrow \frac{p^{2}}{\left(\sqrt{2mE}\right)^{2}} + \frac{q^{2}}{\left(\sqrt{\frac{2E}{m\omega^{2}}}\right)^{2}} = 1
\eenn
Il s'agit donc d'une ellipse dont les demi-axes ont pour longueur $\sqrt{2mE}$ et $\sqrt{\frac{2E}{m\omega^{2}}}$. Sa surface vaut $\pi\sqrt{2mE}\sqrt{\frac{2E}{m\omega^{2}}} = \frac{2\pi E}{\omega}$ et l'invariant adiabatique $I$ d\'{e}fini en (\ref{EQ:49_10}) est \'{e}gal \`{a} :
\be
	I = \frac{E}{\omega} \label{EQ:49_12}
\ee

\section{Variables canoniques}\label{PAR:50}

\subsection{Syst\`{e}me conservatif}\label{PAR:50_1}

Nous consid\'{e}rons maintenant que le param\`{e}tre $\lambda$ introduit au paragraphe (\ref{PAR:49}) est constant et que le syst\`{e}me est ferm\'{e}. Nous allons r\'{e}aliser une transformation canonique, voir le paragraphe (\ref{PAR:45}), des variables $(p,q)$ vers la nouvelle impulsion qui sera $I$, l'invariant adiabatique d\'{e}fini au paragraphe (\ref{PAR:49}). Nous allons d\'{e}montrer que la fonction g\'{e}n\'{e}ratrice de cette transformation canonique est l'action r\'{e}duite $S_{0}$ exprim\'{e}e en fonction de la coordonn\'{e}e $q$ et de la nouvelle impulsion $I$. $S_{0}(q,I,\lambda)$ doit alors se conformer aux relations (\ref{EQ:45_8}) et en particulier \`{a} : $p = \frac{\partial S_{0}(q,I,\lambda)}{\partial q}$.

Par d\'{e}finition de l'action r\'{e}duite pour un syst\`{e}me ferm\'{e}, voir la formule (\ref{EQ:44_4}) :
\be
	S_{0}(q,E,\lambda) = \int{p(q,E,\lambda)\mathrm{d}q} \label{EQ:50_1}
\ee
avec les quantit\'{e}s $E$ et $\lambda$ donn\'{e}es par les propri\'{e}t\'{e}s du syst\`{e}me.

Le syst\`{e}me \'{e}tant ferm\'{e}, la fonction de Hamilton est constante et \'{e}gale \`{a} $E$ le long de la trajectoire. L'invariant adiabatique $I$ est alors fonction de la seule \'{e}nergie. L'action r\'{e}duite $S_{0}$ peut alors tout \`{a} fait s'exprimer sur la forme $S_{0}(q,I,\lambda)$ et par cons\'{e}quent $p = \left(\frac{\partial S_{0}}{\partial q}\right)_{E}$ prise \`{a} $E$ constante. Suivant la remarque pr\'{e}c\'{e}dente, $\left(\frac{\partial S_{0}}{\partial q}\right)_{E}$ co\"{i}ncide avec $\left(\frac{\partial S_{0}}{\partial q}\right)_{I}$ prise \`{a} $I$ constante. Ainsi :
\be
	p = \frac{\partial S_{0}(q,I,\lambda)}{\partial q} \label{EQ:50_2}
\ee
Cela correspond \`{a} la premi\`{e}re des formules de transformation canonique de (\ref{EQ:45_8}). La seconde formule d\'{e}finit une nouvelle coordonn\'{e}e que nous notons $w$ :
\be
	w = \frac{\partial S_{0}(q,I,\lambda)}{\partial I} \label{EQ:50_3}
\ee
$I$ et $w$ sont appel\'{e}s des \emph{variables canoniques} telles que $I$ est la \emph{variable d'action} et $w$, la \emph{variable d'angle} ou encore \emph{variable angulaire}.

Puisque $S_{0}$ ne d\'{e}pend pas explicitement du temps alors la troisi\`{e}me formule dans (\ref{EQ:45_8}) nous permet d'affirmer que la nouvelle fonction de Hamilton $H'$ co\"{i}ncide avec l'ancienne $H$ exprim\'{e}e \`{a} l'aide des nouvelles variables. En d'autres termes, $H'$ repr\'{e}sente l'\'{e}nergie du syst\`{e}me qui est fonction de la variable d'action, soit $E(I)$. Les \'{e}quations de Hamilton pour les nouvelles variables sont :
\be
	\dot{w} = \frac{\partial H'}{\partial I} = \frac{\partial E(I)}{\partial I} = \frac{\mathrm{d}E(I)}{\mathrm{d}I}\text{ et }I = -\frac{\partial H'}{\partial w} = -\frac{\partial E(I)}{\partial w} = 0 \label{EQ:50_4}
\ee
Ainsi, tout comme l'\'{e}nergie du syst\`{e}me ferm\'{e}, la valeur de $I$ est constante par rapport au temps et la variable angulaire est fonction lin\'{e}aire du temps :
\be
	w = \frac{\mathrm{d}E(I)}{\mathrm{d}I}\mathrm{t} + cste = \omega(I)\mathrm{t} + cste \label{EQ:50_5}
\ee
en utilisant la d\'{e}finition (\ref{EQ:49_9}). $w$ constitue la phase des oscillations. Son accroissement pendant une p\'{e}riode de temps vaut :
\bea
	\Delta w & = & w(\mathrm{t} = T) - w(\mathrm{t} = 0) = w(I)T + cste - cste = w(I)\frac{2\pi}{w(I)} \nonumber \\
	\Leftrightarrow \Delta w & = & 2\pi \label{EQ:50_7}
\eea
Ainsi la variation de l'action r\'{e}duite $\Delta S_{0}$ pendant la m\^{e}me dur\'{e}e vaut en utilisant la formule (\ref{EQ:50_1}) avec les nouvelles variables :
\bea
	\Delta S_{0} & = & S_{0}(w,I,\lambda)_{\mathrm{t} = T} - S_{0}(w,I,\lambda)_{\mathrm{t} = 0} = \left(\int{I(E)\mathrm{d}w}\right)_{\mathrm{t} = T} - \left(\int{I(E)\mathrm{d}w}\right)_{\mathrm{t} = 0} = Iw(T)T \nonumber \\
	\Leftrightarrow \Delta S_{0} & = & 2\pi I \label{EQ:50_6}
\eea
L'action r\'{e}duite ne reprend ainsi pas sa valeur initiale mais cro\^{i}t avec le temps.

En exprimant les variables $p$ et $q$, ou toute fonction univoque de ces deux variables, \`{a} l'aide des variables canoniques d'action et d'angle, alors leur valeur respective reste inchang\'{e}e lorsque $w$ varie de $2\pi$ pour une valeur donn\'{e}e de $I$. En d'autres termes, toute fonction univoque de $p$ et $q$ exprim\'{e}e au moyen de variables canoniques est une fonction p\'{e}riodique de la variable angulaire $w$ et de p\'{e}riode $2\pi$.

\subsection{Syst\`{e}me ouvert}\label{PAR:50_2}

Nous \'{e}tudions ici le cas d'un syst\`{e}me ouvert o\`{u} le param\`{e}tre $\lambda$ d\'{e}pend du temps. Nous pouvons reprendre le m\^{e}me raisonnement que dans le paragraphe (\ref{PAR:50_1}) \`{a} l'exception du fait que d\'{e}sormais $S_{0}$ n'est plus d\'{e}finie strictement comme l'action r\'{e}duite puisque l'un de ses param\`{e}tres d\'{e}pend du temps : $S_{0}(q,I,\lambda(\mathrm{t}))$. Toujours fonction g\'{e}n\'{e}ratrice mais d\'{e}sormais fonction du temps, elle implique que la nouvelle fonction de Hamilton ne co\"{i}ncide plus avec l'ancienne mais vaut d\'{e}sormais en reprenant la formule (\ref{EQ:45_8}) :
\benn
	H' = H + \frac{\partial S_{0}}{\partial\mathrm{t}} = E(I,\lambda) + \frac{\partial S_{0}}{\partial q}\frac{\partial q}{\partial\mathrm{t}} + \frac{\partial S_{0}}{\partial I}\frac{\partial I}{\partial\mathrm{t}} + \frac{\partial S_{0}}{\partial\lambda}\frac{\partial\lambda}{\partial\mathrm{t}}
\eenn
or seul le param\`{e}tre $\lambda$ d\'{e}pend explicitement du temps et $\frac{\partial\lambda}{\partial\mathrm{t}} = \frac{\mathrm{d}\lambda}{\mathrm{dt}} = \dot{\lambda}$, donc :
\be
	H' = E(I,\lambda) + \Lambda\dot{\lambda} \label{EQ:50_8}
\ee
en d\'{e}finissant :
\be
	\Lambda = \left(\frac{\partial S_{0}}{\partial\lambda}\right)_{q,I} \label{EQ:50_9}
\ee
Les \'{e}quations de Hamilton s'\'{e}crivent alors :
\be
	\dot{w} = \frac{\partial H'}{\partial I} = \omega(I,\lambda) + \left(\frac{\partial\Lambda}{\partial I}\right)_{w,\lambda}\dot{\lambda} \label{EQ:50_11}
\ee
et en appliquant la d\'{e}finition (\ref{EQ:49_9}), soit $\omega = \left(\frac{\partial E}{\partial I}\right)_{\lambda}$ qui est la fr\'{e}quance des oscillations calcul\'{e}e pour $\lambda$ constante, et comme $\frac{\partial\dot{\lambda}}{\partial I} = 0$ alors :
\be
	\dot{I} = -\frac{\partial H'}{\partial w} = -\frac{\partial E(I,\lambda)}{\partial w} - \left(\frac{\partial\Lambda}{\partial w}\right)_{I,\lambda}\dot{\lambda} = -\left(\frac{\partial\Lambda}{\partial w}\right)_{I,\lambda}\dot{\lambda} \label{EQ:50_10}
\ee

\section{Exactitude de conservation d'invariant adiabatique}\label{PAR:51}

L'\'{e}quation du mouvement (\ref{EQ:50_10}) permet de s'assurer que la variable d'action est un invariant adiabatique. Nous avons vu dans le paragraphe (\ref{PAR:50_1}) que $S_{0}(q,I,\lambda)$ est une fonction non unidorme de la coordonn\'{e}e $q$ : lorsque cette coordonn\'{e}e reprend sa valeur initiale, la valeur de $S_{0}$ est augment\'{e}e d'un multiple entier non nul de la quantit\'{e} $2\pi I$ comme d\'{e}montr\'{e} en (\ref{EQ:50_6}).

De plus, la quantit\'{e} $\Lambda$ d\'{e}finie dans la formule (\ref{EQ:50_9}) est une fonction uniforme en $q$ car la d\'{e}riv\'{e}e de sa d\'{e}finition est r\'{e}alis\'{e}e \`{a} $I$ constant. Donc les accroissements en $I$ qui s'ajoutent \`{a} $S_{0}$ disparaissent dans le calcul de $\Lambda$.

Une des conclusions du paragraphe (\ref{PAR:50_1}) est que toute fonction uniforme exprim\'{e}e \`{a} l'aide des varaibles canoniques est une fonction p\'{e}riodique de $w$ et de p\'{e}riode $2\pi$. Donc la valeur moyenne sur une telle p\'{e}riode de la quantit\'{e} $\frac{\partial\Lambda}{\partial w}$ est nulle. Par extension et en utilisant l'\'{e}quation du mouvement (\ref{EQ:50_10}) et parce que $\lambda$ \'{e}volue lentement avec le temps, i.e. $\overline{\dot{\lambda}} \approx \dot{\lambda}$, alors :
\be
	\overline{\dot{I}} = -\left(\overline{\frac{\partial\Lambda}{\partial w}}\right)_{I}\dot{\lambda} = 0 \label{EQ:51_1}
\ee
Cela d\'{e}montre donc bien que $I$ est un invariant adiabatique.

Les \'{e}quations du mouvement exprim\'{e}es avec les variables canoniques (\ref{EQ:50_10}) et (\ref{EQ:50_11}) permettent d'examiner le probl\`{e}me d'exactitude avec laquelle se conserve l'invariant adiabatique. Pour r\'{e}pondre \`{a} cette question, nous nous pla\c{c}ons dans les hypoth\`{e}ses suivantes :
\benn
	\lambda(\mathrm{t})\underset{\mathrm{t}\rightarrow -\infty}{\longrightarrow}\lambda_{-}\text{, }\lambda(\mathrm{t})\underset{\mathrm{t}\rightarrow +\infty}{\longrightarrow}\lambda_{+}\text{ et }I(\mathrm{t}\rightarrow -\infty) = I_{-}
\eenn
Il nous faut alors d\'{e}terminer l'accroissement $\Delta I = I_{+} - I_{-}$ pour $\mathrm{t}\rightarrow +\infty$.

En utilisant la relation (\ref{EQ:50_10}), nous pouvons \'{e}crire :
\be
	\Delta I = -\int_{-\infty}^{+\infty}{\frac{\partial\Lambda}{\partial w}\dot{\lambda}\mathrm{dt}} \label{EQ:51_2}
\ee
Comme d\'{e}montr\'{e} plus haut dans ce paragraphe, $\Lambda$ est $2\pi$ p\'{e}riodique par rapport \`{a} la variable $w$. En d\'{e}veloppant $\Lambda$ en s\'{e}rie de Fourier, nous \'{e}crivons :
\be
	\Lambda = \sum_{n = -\infty}^{n = +\infty}e^{inw}\Lambda_{n} \label{EQ:51_3}
\ee
Sachant que $\Lambda$ est une fonction aux valeurs r\'{e}elles et que les c{\oe}fficients $\Lambda_{n}$ ne le sont pas alors nous en d\'{e}duisons que $\Lambda_{-n} = \Lambda_{n}^{*}$ qui est le conjugu\'{e} de $\Lambda_{n}$. La d\'{e}riv\'{e}e de $\Lambda$ par rapport \`{a} $w$ sachant que les c{\oe}fficients $\Lambda_{n}$ sont constants est :
\bea
	\frac{\partial\Lambda}{\partial w} & = & \sum_{n=-\infty}^{+\infty}ine^{inw}\Lambda_{n} + \sum_{n=-\infty}^{+\infty}e^{inw}\frac{\partial\Lambda_{n}}{\partial w} = \sum_{n=-\infty}^{+\infty}ine^{inw}\Lambda_{n} \nonumber \\
	& = & \sum_{n=-\infty}^{0}ine^{inw}\Lambda_{n} + \sum_{n=1}^{+\infty}ine^{inw}\Lambda_{n} = \sum_{n=-\infty}^{-1}ine^{inw}\Lambda_{n} + \sum_{n=1}^{+\infty}ine^{inw}\Lambda_{n} \nonumber \\
	& = & -\sum_{n=1}^{+\infty}ine^{-inw}\Lambda_{-n} + \sum_{n=1}^{+\infty}ine^{inw}\Lambda_{n} = \sum_{n=1}^{+\infty}\left(e^{inw}\Lambda_{n} - e^{-inw}\Lambda_{-n}\right)in \nonumber
\eea
Supposons que $\Lambda_{n} = \alpha_{n} + i\beta_{n}$ et nous avons toujours $\Lambda_{-n} = \Lambda_{n}^{*} = \alpha_{n} - i\beta_{n}$, donc en d\'{e}composant les parties r\'{e}elles et imaginaires de l'expression pr\'{e}c\'{e}dente, nous avons :
\bea
	\frac{\partial\Lambda}{\partial w} & = & \sum_{n=1}^{+\infty}in\left[\left(\cos(nw) + i\sin(nw)\right)(\alpha_{n} + i\beta_{n}) - \left(\cos(-nw) + i\sin(-nw)\right)(\alpha_{n} - i\beta_{n})\right] \nonumber \\
	& = & \sum_{n=1}^{+\infty}in\bigg[\alpha_{n}\cos(nw) + i\beta_{n}\cos(nw) + i\alpha_{n}\sin(nw) - \beta_{n}\sin(nw) \nonumber \\
	& & -\left(\alpha_{n}\cos(nw) - i\beta_{n}\cos(nw) - i\alpha_{n}\sin(nw) - \beta_{n}\sin(nw)\right)\bigg] \nonumber \\
	& = & \sum_{n=1}^{+\infty}in\left[2i\beta_{n}\cos(nw) + 2i\alpha_{n}\sin(nw)\right] = -2\sum_{n=1}^{+\infty}n\left[\alpha_{n}\sin(nw) + \beta_{n}\cos(nw)\right] \nonumber
\eea
Or :
\bea
	ie^{inw}\Lambda_{n} & = & i\left(\cos(nw) + i\sin(nw)\right)(\alpha_{n} + i\beta_{n}) \nonumber \\
	& = & i\alpha_{n}\cos(nw) - \beta_{n}\cos(nw) - i\alpha_{n}\sin(nw) - i\beta_{n}\sin(nw) \nonumber
\eea
Donc :
\benn
	\Re\left(ie^{inw}\Lambda_{n}\right) = -\left(\alpha_{n}\sin(nw) + \beta_{n}\cos(nw)\right)
\eenn
Et finalement, nous obtenons :
\be
	\frac{\partial\Lambda}{\partial w} = 2\sum_{n=1}^{+\infty}n\Re\left(ie^{inw}\Lambda_{n}\right) = 2\Re\left(\sum_{n=1}^{+\infty}ine^{inw}\Lambda_{n}\right) \label{EQ:51_4}
\ee
La formule (\ref{EQ:50_11}) qui donne $\dot{w}$ en fonction de $\omega$ et de $\dot{\lambda}$ nous permet d\'{e}crire que si $\dot{\lambda}$ est suffisamment petit alors $\dot{w}$ est positif puisque par d\'{e}finition $\omega$ l'est aussi. Donc $w$ est une fonction croissante avec le temps. Cela permet d'\'{e}crire diff\'{e}remment l'int\'{e}gration (\ref{EQ:51_2}) en ne changeant pas les bornes mais en faisant intervenir $\mathrm{d}w$, soit :
\be
	\Delta I = -\int_{-\infty}^{+\infty}{\frac{\partial\Lambda}{\partial w}\frac{\mathrm{d}\lambda}{\mathrm{dt}}\mathrm{dt}} = -\int_{-\infty}^{+\infty}{\frac{\partial\Lambda}{\partial w}\frac{\mathrm{d}\lambda}{\mathrm{dt}}\frac{\mathrm{dt}}{\mathrm{d}w}\mathrm{d}w} \label{EQ:51_5}
\ee
qui peut se d\'{e}velopper en utilisant (\ref{EQ:51_4}) ainsi :
\benn
	\Delta I = -2\Re\left[\sum_{n=1}^{+\infty}\int_{-\infty}^{\infty}{ine^{inw}\Lambda_{n}\frac{\mathrm{d}\lambda}{\mathrm{dt}}\frac{\mathrm{dt}}{\mathrm{d}w}\mathrm{d}w}\right]
\eenn

\begin{figure}[htb!]
        \begin{center}
                \includegraphics[width=15cm]{chapter_07_paragraph_51}
		\caption{Visualisation du chemin d'int\'{e}gration parcouru par la variable $w$}\label{FIG:51_1}
        \end{center}
\end{figure}

L'int\'{e}grale se calcule sur le chemin de l'axe r\'{e}el pour $w$. En prenant $w$ comme une variable complexe et en supposant qu'aucune int\'{e}gration ne poss\`{e}de de points singuliers pour la partie r\'{e}elle de $w$, nous d\'{e}pl\c{c}ons le chmin d'int\'{e}gration dans la partie imaginaire telle sch\'{e}matis\'{e}e sur la figure (\ref{FIG:51_1}). Le chemin d'int\'{e}gration rencontre alors les points singuliers de $w$ et s'y accroche en quelque sorte.

Pour $w \in \mathbb{C}$ tel que $w = a + ib$ alors :
\benn
	ine^{inw} = ine^{ina}e^{-nb} = ne^{-nb}\left[icos(na) - sin(na)\right]
\eenn
donc :
\benn
	\Re\left(ine^{inw}\right) = -nsin(na)e^{-nb} \propto e^{-\Im(w)}
\eenn
Soit $w_{0}$ le point singulier le plus proche de l'axe r\'{e}el, et qui a par cons\'{e}quent la partie imaginaire la plus petite de l'ensemble des points singuliers.

La contribution de chacun des termes de la relation (\ref{EQ:51_4}) \`{a} l'int\'{e}grale (\ref{EQ:51_5}) est alors maximale pour $w_{0}$ dont la partie imaginaire est minimale. Ainsi pour $\Delta I$, la contribution maximale est pour $n = 1$ telle que :
\be
	\Delta I \propto e^{-\Im(w_{0})} \label{EQ:51_6}
\ee
Soit $\mathrm{t}_{0}$, le << temps >>, en nombre complexe, tel que $w(\mathrm{t}_{0}) = w_{0}$ et qui correspond au temps caract\'{e}ristique de variation des param\`{e}tres du syst\`{e}me. Nous d\'{e}signons cet intervalle de temps par $\tau$. Ainsi, en prenant l'\'{e}quation de Hamilton (\ref{EQ:50_11}), nous avons :
\be
	\Im(w_{0}) \approx \omega\tau \approx \frac{\tau}{T} \label{EQ:51_7}
\ee
avec $T$ qui est la p\'{e}riode du mouvement d\'{e}finie en (\ref{EQ:49_8}). Par d\'{e}finition, les param\`{e}tres du syst\`{e}me \'{e}voluent lentement donc $\tau \gg T$. Par cons\'{e}quent, $\Delta I = I_{+} - I_{-} \propto e^{-\Im(w_{0})}$ d\'{e}cro\^{i}t exponentiellement avec la vitesse de variation des param\`{e}tres du syst\`{e}me.

Pour calculer $w_{0}$ en premi\`{e}re approximation par rapport \`{a} $\frac{T}{\tau} = \left(\frac{\tau}{T}\right)^{-1}$, nous partons de l'\'{e}quation de Hamilton (\ref{EQ:50_11}) dans laquelle nous n\'{e}gligeons le terme en $\dot{\lambda}$, soit :
\be
	\frac{\mathrm{d}w}{\mathrm{dt}} = \omega(I,\lambda(\mathrm{t})) \label{EQ:51_8}
\ee
En supposant que pour le calcul de $w_{0}$, la valeur de $I$ reste constante, alors :
\be
	w_{0} = \int^{\mathrm{t}_{0}}{\omega(I,\lambda(\mathrm{t}))\mathrm{dt}} \label{EQ:51_9}
\ee
En ins\'{e}rant l'\'{e}quation (\ref{EQ:51_8}) dans la formulation de $\Delta I$ en (\ref{EQ:51_5}), nous avons :
\benn
	\Delta I = -2\int_{-\infty}^{+\infty}{\Re\left[\sum_{n=1}^{+\infty}ine^{inw}\right]\dot{\lambda}\frac{\mathrm{d}\omega}{\omega(I,\lambda)}}
\eenn
qui permet, en prenant le terme de la s\'{e}rie pour $n = 1$, s'\'{e}crire :
\be
	\Delta I \propto \Re\left[\int_{-\infty}^{+\infty}{ie^{inw}\dot{\lambda}\frac{\mathrm{d}\omega}{\omega(I,\lambda)}}\right] \label{EQ:51_10}
\ee

\section{Mouvement quasi-pr\'{e}riodique}\label{PAR:52}

\subsection{G\'{e}n\'{e}ralit\'{e}s}\label{PAR:52_1}

Nous consid\'{e}rons un syst\`{e}me ferm\'{e} \`{a} plusieurs degr\'{e}s de libert\'{e} qui effectue un mouvement fini pour toutes ses coordonn\'{e}es. Nous supposons alors que le probl\`{e}me admette une s\'{e}paration des variables compl\`{e}te selon la m\'{e}thode de Hamilton-Jacobi comme d\'{e}crite dans le paragraphe (\ref{PAR:48}). Ainsi la d\'{e}finition (\ref{EQ:48_6}) indique que pour un choix convenable de coordonn\'{e}es, l'action r\'{e}duite du syst\`{e}me est une somme de fonctions ne d\'{e}pendant chacune que d'une seule coordonn\'{e}e :
\be
	S_{0} = \sum_{i}S_{i}(q_{i}) \label{EQ:52_1}
\ee
Par d\'{e}finition et en utilisant la formule (\ref{EQ:43_3}), $p_{i} = \frac{\partial S_{i}}{\partial q_{i}} = \frac{\partial S_{0}}{\partial q_{i}}$ et comme nous avons la relation (\ref{EQ:44_4}) alors $S_{0} = \int{\sum_{i}p_{i}\mathrm{d}q_{i}}$. Cela nous donne finalement :
\be
	S_{i} = \int{p_{i}\mathrm{d}q_{i}} \label{EQ:52_2}
\ee
Par hypoth\`{e}se initiale, le mouvement est fini donc chaque coordonn\'{e}e prend des valeurs born\'{e}es. Pour chaque coordonn\'{e}e, lorsqu'elle revient \`{a} sont point d'origine apr\`{e}s avoir atteint ses bornes minimales et maximales alors la relation (\ref{EQ:50_6}) nous $\Delta S_{i} = 2\pi I_{i}$ et par voie de cons\'{e}quence :
\be
	\Delta S_{0} = \sum_{i}\Delta_{i} = 2\pi\sum_{i}I_{i} \label{EQ:52_3}
\ee
o\`{u} $I_{i}$ est l'int\'{e}grale :
\be
	I_{i} = \frac{1}{2\pi}\oint{p_{i}\mathrm{d}q_{i}} \label{EQ:52_4}
\ee
conform\'{e}mement \`{a} la d\'{e}finition (\ref{EQ:49_7}).

Nous r\'{e}alisons d\'{e}sormais la transformation canonique en \'{e}tendant celle op\'{e}r\'{e}e au paragraphe (\ref{PAR:50}) pour un unique degr\'{e} de libert\'{e}. Nous obtenons les variables d'action $I_{i}$ et les variables d'angle :
\be
	w_{i} = \frac{\partial S_{0}(q,I)}{\partial I_{i}} = \sum_{j}\frac{\partial S_{j}(q_{j},I)}{\partial I_{i}} \label{EQ:52_5}
\ee
Nous avons de nouveau l'action exprim\'{e}e en fonction des coordonn\'{e}es et des variables d'action $I_{i}$ qui jouent le r\^{o}le de fonction g\'{e}n\'{e}ratrice de la transformation. Les \'{e}quations du mouvement sont alors pour chaque degr\'{e} de libert\'{e} s\'{e}par\'{e} et en adaptant les relations (\ref{EQ:50_4}) :
\be
	\dot{I}_{i} = 0\text{ et }\dot{w}_{i} = \frac{\partial E(I)}{\partial I_{i}} \label{EQ:52_6}
\ee
qui nous donnent conform\'{e}ment \`{a} la relation (\ref{EQ:50_5}) :
\bea
	I_{i} & = & cste \nonumber \\
	w_{i} & = & \frac{\partial E(I)}{\partial I_{i}}\mathrm{t} + cste \label{EQ:52_7}
\eea
De mani\`{e}re \'{e}vidente et analogue \`{a} (\ref{EQ:50_7}), nous avons pour une variation compl\`{e}te de $q_{i}$ :
\be
	\Delta w_{i} = 2\pi \label{EQ:52_8}
\ee
Les grandeurs $w_{i}(q,I)$ peuvent varier d'un multiple de $2\pi$ lorsque les coordonn\'{e}es reprennent leur valeur initiale.

Par d\'{e}finition (\ref{EQ:52_4}), les quantit\'{e}s $I_{i}$ sont des fonctions univoques de $p$ et $q$ si nous les exprimons avec ces variables. En portant $I_{i}(p,q)$ dans l'expression de $w_{i}(q,I)$ alors $w_{i}(p,q)$ est obtenue et elle varie d'un multiple de $2\pi$ (ou de 0) lorsqu'elle parcourt une courbe ferm\'{e}e quelconque dans l'espace des phases.
Ainsi toute fonction univoque $F(p,q)$ de l'\'{e}tat du syst\`{e}me que nous exprimons \`{a} l'aide des variables canoniques est fonction p\'{e}riodique des variables angulaires de p\'{e}riode $2\pi$ pour chacune d'elles. Nous pouvons donc la d\'{e}velopper en s\'{e}rie de Fourier telle que :
\benn
	F = \sum_{n_{1}}\cdots\sum_{n_{s}}A_{n_{1}\cdots n_{s}}\exp{\left(i\mathrm{t}(n_{1}w_{1} + \cdots + n_{s}w_{s})\right)}
\eenn
avec $(n_{1},\cdots,n_{s}) \in \mathbb{Z}^{s}$. Nous ins\'{e}rons la d\'{e}finition (\ref{EQ:52_7}) des variables d'angle et obtenons :
\be
	F = \sum_{n_{1}}\cdots\sum_{n_{s}}A_{n_{1}\cdots n_{s}}\exp{\left(i\mathrm{t}\left(n_{1}\frac{\partial E}{\partial I_{1}} + \cdots + n_{s}\frac{\partial E}{\partial I_{s}}\right)\right)} \label{EQ:52_9}
\ee
Ainsi chaque terme de la somme pr\'{e}c\'{e}dente est une fonctoin p\'{e}riodique du temps avec pour fr\'{e}quence :
\be
	\sum_{n_{i}=1}^{s}n_{i}\omega_{i} \label{EQ:52_10}
\ee
qui est la somme de multiples entiers de fr\'{e}quences fondamentales :
\be
	\omega_{i} = \frac{\partial E}{\partial I_{i}} \label{EQ:52_11}
\ee
En g\'{e}n\'{e}ral, ces fr\'{e}quences ne sont pas des multiples entiers de l'une d'entre elles. L'ensemble n'est donc pas strictement p\'{e}riodique, que ce soit pour les coordonn\'{e}es $q$ que pour les impulsions $p$. Le syst\`{e}me, apr\`{e}s \^{e}tre pass\'{e} par un \'{e}tat, ne repassera pas forc\`{e}ment par ce m\^{e}me \'{e}tat apr\`{e}s un temps fini. Toutefois, apr\`{e}s un temps suffisamment long, le syst\`{e}me passera aussi pr\`{e}s que souhait\'{e} de cet \'{e}tat. Un tel mouvement est dit \emph{quasi-p\'{e}riodique}.

Si deux ou plus des fr\'{e}quences fondamentales $\omega_{i}$ sont commensurables, i.e. telles que leur rapport entre elles est un nombre rationnel, alors il s'agit de cas de \emph{d\'{e}g\'{e}n\'{e}rescence}. Si toutes les fr\'{e}quences sont commensurables alors le mouvement du syst\`{e}me est dit \emph{compl\`{e}tement d\'{e}g\'{e}n\'{e}r\'{e}}. Dans ce dernier cas, le mouvement est rigoureusement p\'{e}riodique et la trajectoire de toutes les particules est ferm\'{e}e.

\subsection{Propri\'{e}t\'{e}s particuli\`{e}res}\label{PAR:52_2}

\subsubsection{Diminution des grandeurs ind\'{e}pendantes}\label{PAR:52_2_1}

Soient $\omega_{1}$ et $\omega_{2}$ deux fr\'{e}quences commensurables. Elles sont donc li\'{e}es par une relation telle que, avec $n_{1}$ et $n_{2}$ deux entiers non nuls, $n_{1}\omega_{1} = n_{2}\omega_{2}$. Cela nous donne par d\'{e}finition :
\be
	n_{1}\frac{\partial E}{\partial I_{1}} = n_{2}\frac{\partial E}{\partial I_{2}} \label{EQ:52_12}
\ee
Cette relation nous donne $n_{1}\partial I_{2} = n_{2}\partial I_{1}$. Ainsi les variables d'action $I_{1}$ et $I_{2}$ n'apparaissent dans la formulation de l'\'{e}nergie $E$ que sous la forme $n_{1}I_{2} + n_{2}I_{1}$. De fait, cela diminue le nombre de grandeurs ind\'{e}pendantes $\{I_{i}\}$.

\subsubsection{Cons\'{e}quences sur les int\'{e}grales premi\`{e}res}\label{PAR:52_2_2}

Au paragraphe (\ref{PAR:6}), nous avons vu qu'il y a $2s-1$ int\'{e}grales du mouvement pour un syst\`{e}me ferm\'{e} \`{a} $s$ degr\'{e}s de libert\'{e} et non d\'{e}g\'{e}n\'{e}r\'{e}. Les relations $\dot{I}_{i} = 0$ des \'{e}quations du mouvement (\ref{EQ:52_6}) nous donne d\'{e}j\`{a} $s$ int\'{e}grales uniformes, i.e. toutes les quantit\'{e}s $I_{i}$ sont constantes lors du mouvement. Les \'{e}quations du mouvement pour les variables angulaires donnent $\dot{w}_{i} = \frac{\partial E}{\partial I_{i}}$ donc pour $i \ne j$, nous avons :
\benn
	\dot{w}_{i}\frac{\partial E}{\partial I_{j}} - \dot{w}_{j}\frac{\partial E}{\partial I_{i}} = 0 \Leftrightarrow \frac{\mathrm{d}}{\mathrm{dt}}\left(w_{i}\frac{\partial E}{\partial I_{j}} - w_{j}\frac{\partial E}{\partial I_{i}}\right) = 0
\eenn
Donc les $s - 1$ quantit\'{e}s :
\be
	\forall i \ne j\text{, }w_{i}\frac{\partial E}{\partial I_{j}} - w_{j}\frac{\partial E}{\partial I_{i}} \label{EQ:52_13}
\ee
sont des int\'{e}grales mais elles ne sont pas uniformes car les cariables angulaires ne le sont pas au regard de l'\'{e}galit\'{e} (\ref{EQ:52_8}).

Dans le cas de la d\'{e}g\'{e}n\'{e}rescence abord\'{e}e au paragraphe (\ref{PAR:52_2_1}), nous \'{e}tudions la d\'{e}riv\'{e}e totale par rapport au temps de la quantit\'{e} $w_{1}n_{2} - w_{2}n_{1}$, soit :
\benn
	\frac{\mathrm{d}}{\mathrm{dt}}(w_{1}n_{2} - w_{2}n_{1}) = \frac{\partial E}{\partial I_{1}}n_{2} - \frac{\partial E}{\partial I_{2}}n_{1} = \frac{\partial E}{\partial I_{2}}\left(\frac{n_{2}^{2}}{n_{1}} - n_{1}\right) \propto \dot{w}_{2}
\eenn
en utilisant la relation (\ref{EQ:52_12}). Nous obtenons alors que la quantit\'{e} :
\be
	w_{1}n_{2} - w_{2}n_{1} \label{EQ:52_14}
\ee
est proportionnelle \`{a} $\dot{w}_{2}$. L'int\'{e}grale (\ref{EQ:52_14}) n'est pas uniforme mais comme la non-uniformit\'{e} de $w_{2}$ se r\'{e}duit \`{a} un multiple de $2\pi$, voir (\ref{EQ:52_8}), alors il suffit de choisir une fonction trigonom\'{e}trique ajust\'{e}e pour que, \`{a} partir de (\ref{EQ:52_14}), nous ayons une int\'{e}grale premi\`{e}re uniforme.
